\documentclass[11pt, oneside]{article}   	% use "amsart" instead of "article" for AMSLaTeX format
\usepackage[width=16.5cm, top=3cm, bottom=3cm]{geometry}                		% See geometry.pdf to learn the layout options. There are lots.
\geometry{a4paper}                   		% ... or a4paper or a5paper or ... 
%\geometry{landscape}                		% Activate for rotated page geometry
\usepackage[parfill]{parskip}    		% Activate to begin paragraphs with an empty line rather than an indent
\usepackage{graphicx}				% Use pdf, png, jpg, or eps§ with pdflatex; use eps in DVI mode
								% TeX will automatically convert eps --> pdf in pdflatex		
\usepackage{amssymb}

%SetFonts

%SetFonts
\usepackage[english]{babel}
\usepackage{amsmath,amssymb,stmaryrd,enumerate,mathrsfs,bbm,xcolor,latexsym,theorem}

%%%%%%%%%% Start TeXmacs macros
\newcommand{\Iota}{\mathrm{I}}
\newcommand{\assign}{:=}
\newcommand{\mathD}{\mathrm{D}}
\newcommand{\mathd}{\mathrm{d}}
\newcommand{\nocomma}{}
\newcommand{\tmaffiliation}[1]{\\ #1}
\newcommand{\tmcolor}[2]{{\color{#1}{#2}}}
\newcommand{\tmop}[1]{\ensuremath{\operatorname{#1}}}
\newcommand{\tmtextbf}[1]{{\bfseries{#1}}}
\newcommand{\tmtextit}[1]{{\itshape{#1}}}
\newenvironment{enumeratealpha}{\begin{enumerate}[a{\textup{)}}] }{\end{enumerate}}
\newenvironment{proof}{\noindent\textbf{Proof\ }}{\hspace*{\fill}$\Box$\medskip}
\newtheorem{definition}{Definition}
\newtheorem{lemma}{Lemma}
\newtheorem{notation}{Notation}
\newtheorem{proposition}{Proposition}
{\theorembodyfont{\rmfamily}\newtheorem{remark}{Remark}}
\newtheorem{theorem}{Theorem}
\newcommand{\tmkeywords}{\textbf{Keywords:} }
%%%%%%%%%% End TeXmacs macros

\def\Xint#1{\mathchoice
{\XXint\displaystyle\textstyle{#1}}%
{\XXint\textstyle\scriptstyle{#1}}%
{\XXint\scriptstyle\scriptscriptstyle{#1}}%
{\XXint\scriptscriptstyle\scriptscriptstyle{#1}}%
\!\int}
\def\XXint#1#2#3{{\setbox0=\hbox{$#1{#2#3}{\int}$ }
\vcenter{\hbox{$#2#3$ }}\kern-.6\wd0}}
\def\ddashint{\Xint=}
\def\dashint{\Xint-}


\newcommand{\bint}{\dashint}
\newcommand{\dint}{\not{d}}
\newcommand{\C}{\mathscr{C}}
\newcommand{\CC}{\mathscr{C}}
\newcommand{\CCC}{\mathscr{C}}
\newcommand{\VV}{\mathscr{C}}
\newcommand{\CD}{\mathscr{D}}
\newcommand{\CF}{\mathscr{F}}
\newcommand{\CB}{\mathscr{B}}
\newcommand{\CA}{\mathscr{A}}
\newcommand{\CL}{\mathscr{L}}
\newcommand{\CQ}{\mathscr{Q}}
\newcommand{\CS}{\mathscr{S}}
\newcommand{\CP}{\mathscr{P}}
\newcommand{\CX}{\mathscr{X}}
\newcommand{\CZ}{\mathscr{Z}}
\newcommand{\addref}{\tmcolor{red}{[add reference]}}
%


\begin{document}

\title{Variational approach to Euclidean QFT}

\author{
  N. Barashkov
  \tmaffiliation{IAM University of Bonn, Germany}
  \and
  M. Gubinelli
  \tmaffiliation{IAM \& HCM University of Bonn, Germany}
}

\maketitle

\begin{abstract}
  We propose a new approach to the construction of the $\Phi^4_{3}$
  Euclidean field theory on a periodic domain. The key ingredient is a
  variational formula for the partition function involving a stochastic
  control problem along the flow of a scale regularization parameter.
  Techniques originating from the analysis of singular SPDEs give uniform
  bounds allowing to prove the boundedness of the variational functional as
  the cutoff is removed. The control problem reveals itself as a rigorous and
  efficient version of Wilson's continous renormalization group.
\end{abstract}

\tmkeywords{Euclidean quantum field theory, Bou{\'e}--Dupuis formula,
renormalization group, paracontrolled calculus.}

\section{Introduction }

The $\Phi^4_d$ Euclidean quantum field theory on the $d$ dimensional torus
$\Lambda =\mathbbm{T}^d = (\mathbbm{R}/ (2 \pi \mathbbm{Z}))^d$ is the
probability measure $\nu$ obtained as the weak limit for $T \rightarrow
\infty$ of the family of measures $(\nu_T)_T$ given by
\begin{equation}
  \nu_T (\mathd \phi) = \CZ^{- 1}_T \exp \left[ - \lambda \int_{\Lambda}
  [\phi_T (\xi)^4 - a_T \phi_T (\xi)^2] \mathd \xi \right] \mu (\mathd \phi)
  \label{eq:gibbs} .
\end{equation}
Here $\Delta$ is the Laplacian on $\Lambda$, $\mu$ is the centered Gaussian
measure with covariance $(1 - \Delta)^{- 1}$, $\CZ_T$ is a normalization
factor, $a_T$ a given constant and $\phi_T = \rho_T \ast \phi$ with $\rho_T$
some appropriate smooth and compactly supported cutoff function such that
$\rho_T \rightarrow \delta$ as $T \rightarrow \infty$. The measures $\mu,
\nu_T$ are realized as probability measures on $\CS' (\Lambda)$, the space of
distributions on $\Lambda$. It is well known that they are supported on the
Sobolev space $W^{(2 - d) / 2 - \varepsilon, p} (\Lambda)$ for all small
$\varepsilon > 0$ and $p \geqslant 2$. The existence of the limit $\nu$ is
conditioned on the choice of an appropriate sequence of
\tmtextit{renormalization constants} $(a_T)_T$.

\

The aim of this paper is to present a proof of convergence using stochastic
calculus along the direction given by the continuous scale cutoff $T$ together
with a variational formula for the partition function $\CZ_T$ and for the
generating function of the measure $\nu_T$. Our work can be seen as an attempt
to pursue Wilson's~{\cite{wilson_renormalization_1983}} and
Polchinski's~{\cite{polchinski_renormalization_1984}} continuous
renormalization group (RG) method. This method has been made rigorous by
Brydges, Slade et
al.~{\cite{brydges_functional_1993,brydges_short_1995,brydges_lectures_2009}}
and as such witnesses a lot of progress and
successes~{\cite{brydges_renormalisation_2014_3,brydges_renormalisation_2014_2,bauerschmidt_renormalisation_2014,brydges_renormalisation_2014_1,brydges_renormalisation_2014,brydges_short_1995}}.
The key idea is the nonperturbative study of a certain infinite dimensional
Hamilton--Jacobi--Bellman equation~{\cite{brydges_mayer_1987}} describing the
effective, scale dependent, action of the theory. Here we take the stand of
trying to avoid the heavy analysis involved by the direct study of the PDE by
going to the equivalent stochastic control formulation, well established and
understood in finite dimensions~{\cite{fleming_controlled_2005}}. The time
parameter of the evolution corresponds to an increasing amount of small scale
fluctuations of the Euclidean field and our main tool is a variational
representation formula, introduced by Bou{\'e} and
Dupuis~{\cite{boue_variational_1998}}, for the logarithm of the partition
function interpreted as the value function of the control problem. See also
the related papers of {\"U}st{\"u}nel~{\cite{ustunel_variational_2014}} and
Zhang~{\cite{zhang_variational_2009}} where extensions and further results on
the variational formula are obtained. This variational formula has been used
by Lehec~{\cite{lehec_representation_2013}} to reprove some Gaussian
functional inequalities. In this representation we can avoid the analysis of
an infinite dimensional second order operator and concentrate more on pathwise
properties of the Euclidean interacting fields. We are able to leverage
techniques developed for singular SPDEs, in particular the paracontrolled
calculus developed in~{\cite{gubinelli_paracontrolled_2015}}, to perform the
renormalization of various non-linear quantities and show uniform bounds in
the $T \rightarrow \infty$ limit.

\

The main result of the paper is the following bound.

\begin{theorem}
  \label{th:main-3d}Let $d = 3$. There exist renormalization constants
  $\alpha_T, \delta_T$, which depend polynomially on $\lambda$, and functions
  $E_+ (\lambda), E_- (\lambda)$ not depending on $\Lambda$, such that
  \[ \lim_{\lambda \rightarrow 0} \frac{E_{+, -} (\lambda)}{\lambda^3} = 0,
  \]
  and for any $\lambda > 0$
  \[ E_- (\lambda) \leqslant - \frac{1}{| \Lambda |} \log \int \exp \left[ -
     \lambda \int_{\Lambda} [\phi_T (\xi)^4 - \alpha_T \phi_T (\xi)^2 -
     \delta_T] \mathd \xi \right] \mu (\mathd \phi) \leqslant E_+ (\lambda) .
  \]
\end{theorem}

A similar statement for $d = 2$ will be proven below in order to introduce
some of the ideas on which the $d = 3$ proof is based. The method of proof
give also boundedness of the moment generating function of $\phi_T$ uniformly
in $T$ allowing to prove convergence in law, at least along subsequences
$(T_n)_n$, and some additional informations on the pathwise behaviour of the
interacting fields. We do not pursue further these developments since it is
also clear from the methods we introduce that the full convergence of the
sequence of measures $(\nu_T)_T$ and of the corresponding variational problems
can be obtained using arguments of $\Gamma$--convergence which we leave to a
further publication.

\

The construction of the $\Phi^4_{2, 3}$ measure in finite volume is basic
problem of constructive quantum field theory to which much literature has been
devoted. Is not our aim to provide here a comprehensive review of this
literature.The $d = 2$ case is much simpler than $d = 3$ and a large amount of
work has been devoted to this situation. As far as the $d = 3$ case is
concerned, let us just mention some of the results that, to different extent,
prove the existence of the limit as the ultraviolet (small scale)
regularization is removed. After the early fundational work on Glimm and
Jaffe~{\cite{glimm_boson_1968,glimm_positivity_1973}} which was partly
performed in the Hamiltonian formalism, all the subsequent research has been
formulated in the Euclidean setting: i.e. as the problem of construction and
study of the probability measures $\nu$ on a space of distributions.
Feldman~{\cite{feldman_lambda_1974}}, Park~{\cite{park_lambda_1975}}, Benfatto
et al.~{\cite{benfatto_ultraviolet_1980}}, Magn{\'e}n and
Seneor~{\cite{magnen_infinite_1976}} and finally Brydges et
al.~{\cite{brydges_new_1983}} obtained the main results we are aware of.
Recent advances in the analysis of singular SPDEs put forward by the invention
of regularity structures by M.~Hairer~{\cite{hairer_theory_2014}} and related
approaches~{\cite{gubinelli_paracontrolled_2015,catellier_paracontrolled_2013,otto_quasilinear_2016}}
or even RG--inspired ones~{\cite{kupiainen_renormalization_2016}}, have
allowed to pursue the stochastic quantisation program to a point where now can
be used to prove directly existence of the finite volume $\Phi^4_3$ measure in
two different ways~{\cite{mourrat_dynamic_2016,albeverio_invariant_2017}}.
Some results on the existince of the infinite volume dynamics have been
obtained recently~{\cite{gubinelli_global_2018}}. For a recent overview of the
status of the constructive program wrt. the analysis of the $\Phi^4_{2, 3}$
models the reader can consult the introduction
to~{\cite{albeverio_invariant_2017}}.

\

The paper is organized as follows. In Section~\ref{sec:rg-flow} we set up our
main tool, the Bou{\'e}--Dupuis variational formula of
Theorem~\ref{th:variational}. Then, as a warmup exercise, we use the formula
to show existence of the $\Phi^4_2$ measure in Section~\ref{sec:two-d}. We
pass then to the more involved situation in three dimensions in
Section~\ref{sec:three-d} and give the proof of Theorem~\ref{th:main-3d}. In
Section~\ref{sec:stochastic} we collect the proof of renormalizability of the
stochastic terms and related probabilistic estimates. Finally the appendix
contains background material on functional spaces and paraproducts and also
some auxiliary results. The reader is invited to skim over the appendix in
order to be introduced to the notations used in the paper.

\medskip
{\noindent}\tmtextbf{Acknowledgments . }The authors are partially supported
by the German DFG via CRC 1060.{\hspace*{\fill}}{\medskip}

\section{The RG flow as a stochastic control problem}\label{sec:rg-flow}

We begin by constructing a probability space $\mathbbm{P}$ endowed with a
process $$Y \in C (\mathbbm{R}_+, W^{(2 - d) / 2 - \varepsilon, p}
(\Lambda)),$$ such that $\tmop{Law}_{\mathbbm{P}} (Y_{\infty}) = \nu$ and
$\tmop{Law}_{\mu} (\phi_T) = \tmop{Law}_{\mathbbm{P}} (Y_T)$.

\

Let us fix some notations and objects. Let $\langle a \rangle = (1 + a^2)^{1
/ 2}$. Denote with $\CS (\Lambda)$ the space of Schwartz functions on
$\Lambda$ and with $\CS' (\Lambda)$ the dual space of distributions. The
notation $\hat{f}$ or $\CF f$ stands for the space Fourier transform of $f$. \
In order to easily keep track of the volume dependence of various objects we
normalize Lebesgue measure on $\Lambda$ to have unit mass. We denote the
normalized integral and measure by
\[ \bint f \assign \frac{1}{| \Lambda |} \int_{\Lambda} f, \quad \dint x =
   \frac{1}{| \Lambda |} \mathd x. \]
Norms in all the related functional spaces (Lebesgue, Sobolev and Besov
spaces) are understood similarly normalized unless \ stated otherwise. The
various contants appearing in the estimates will be understood independent of
$\Lambda$



Fix $\alpha < - d / 2$ and let $\Omega \assign C (\mathbbm{R}_+ ; H^{-
\alpha})$, $(X_t)_t$ the canonical process on $\Omega$ and $\CB$ the Borel
$\sigma$--algebra of $\Omega$. On $\left( \Omega, \CB \right)$ consider the
probability measure $\mathbbm{P}$ which makes the canonical process $X$ a
cylindrical Brownian motion in $L^2 (\Lambda)$. In the following $\mathbbm{E}$
without any qualifiers will denote expectations wrt. $\mathbbm{P}$ and
$\mathbbm{E}_{\mathbbm{Q}}$ will denote expectations wrt. some other measure
$\mathbbm{Q}$ on $\left( \Omega, \CB \right)$. On the measure space $\left(
\Omega, \CB, \mathbbm{P} \right)$ there exists a collection $(B_t^n)_{n \in
\mathbbm{Z}^d}$ of complex (2-dimensional) Brownian motions, such that
$\overline{B^n_t} = B^{- n}_t$, \ $B^n_t, B^m_t$ independent for $m \neq \pm
n$ and $X_t = \sum_{n \in \mathbbm{Z}^d} e^{2 \pi i \langle n, \cdot \rangle}
B^n_t$. Note that $X$ has a.s. trajectories in $C (\mathbbm{R}_+, W^{- d / 2 -
\varepsilon, p} (\Lambda))$ for any $\varepsilon > 0$ and $p \geqslant 2$ by
standard arguments.

Fix some $\rho \in C_c^{\infty} (\mathbbm{R}_+, \mathbbm{R}_+)$ such that
$\rho (0) = 1$. Let $\rho_t (x) = \rho (x / t)$ and
\[ \sigma_t (x) \assign (\dot{\rho}_t (x))^{1 / 2} = (- \rho' (x / t) (x /
   t^2))^{1 / 2} = (- (x / t) \rho' (x / t))^{1 / 2} / t^{1 / 2}, \]
where $\dot{\rho}_t$ is the partial derivative of $\rho_t$ with respect to
$t$. Consider the process $Y_t$ defined by
\begin{equation}
  Y_t \assign \sum_{n \in \mathbbm{Z}^d} \int_0^t \frac{\sigma_s (n)}{\langle
  n \rangle} e^{2 \pi i \langle n, \cdot \rangle} \tmop{dB}^n_s .
  \label{eq:rep-Y}
\end{equation}
It is a centered Gaussian process with covariance
\begin{eqnarray*}
  \mathbbm{E} [\langle Y_t, \varphi \rangle \langle Y_s, \psi \rangle] & = &
  \sum_{n \in \mathbbm{Z}^d} \sum_{m \in \mathbbm{Z}^d} \mathbbm{E} \left[
  \int_0^t \frac{\sigma_u (n)}{\langle n \rangle} \tmop{dB}^n_u  \hat{\varphi}
  (n) \overline{\int_0^s \frac{\sigma_u (m)}{\langle m \rangle} \tmop{dB}^m_s 
  \hat{\psi} (m)} \right]\\
  & = & \sum_{n \in \mathbbm{Z}^d} \frac{\rho_{\min (s, t)} (n)}{\langle n
  \rangle^2} \hat{\varphi} (n) \overline{\hat{\psi} (n)},
\end{eqnarray*}
for any $\varphi, \psi \in \CS (\Lambda)$, by Fubini's theorem and Ito
isometry. Then for any $\varphi \nocomma, \psi \in L^2 (\Lambda)$ we have
$\lim_{t \rightarrow \infty} \mathbbm{E} [\langle Y_t, \varphi \rangle \langle
Y_t, \psi \rangle] = \sum_{n \in \mathbbm{Z}^d} \langle n \rangle^{- 2}
\hat{\varphi} (n) \overline{\hat{\psi} (n)}$, by dominated convergence.

Note that up to any finite time $T$ the r.v. $\mathcal{\CF} Y_T$ has finite
support and the stopped process $Y^T_t = Y_{t \wedge T}$ for any fixed $T >
0$, is in $C (\mathbbm{R}_+, W^{k, 2} (\Lambda))$ for any $k \in \mathbbm{N}$.
Furthermore $(Y^T_t)_t$ only depends on a finite subset of the Browninan
motions $(B^n)_n$.

We will usually write $g (\mathD)$ for the Fourier multiplier operator with
symbol $g$. With this convention we can compactly denote
\begin{equation}
  Y_t = \int_0^t J_s \mathd X_s \label{eq:def-Y}
\end{equation}
where $J_s \assign \langle D \rangle^{- 1} \sigma_s (D)$. We observe that
$Y_t$ has a distribution given by the pushforward $(\rho_t (\mathD))_{\ast}
\mu$ of $\mu$ through $\rho_t (D)$. We write the measure $\nu_T$
in~(\ref{eq:gibbs}) in terms of expectations over $\mathbbm{P}$ as
\begin{equation}
  \int f (\phi) \nu_T (\mathd \phi) = \frac{\mathbbm{E} [f (Y_T) e^{- V_T
  (Y_T)}]}{\CZ_T}, \label{eq:equiv}
\end{equation}
for any bounded measurable $f : \CS' (\Lambda) \rightarrow \mathbbm{R}$ where
\[ V_T (Y_T) \assign \int_{\Lambda} (\lambda | Y_T (\xi) |^4 - a_T  | Y_T
   (\xi) |^2 - b_T) \mathd \xi, \qquad \CZ_T \assign \mathbbm{E} [e^{- V_T
   (Y_T)}] . \]
The sequence $(b_T)_T$ will be chosen appropriately below in order to decouple
the behaviour of the numerator from that of the denominator.

The first problem we address is that of giving upper and lower bounds to the
quantity $\CZ_T$ appearing in~(\ref{eq:equiv}). Note that, for fixed $T$, the
polynomial appearing in the expression for $V_T (Y_T)$ is bounded below if
$\lambda > 0$, in this case $\CZ_T$ is well defined and also bounded away from
zero (this follows easily from Jensen's inequality). However as $T \rightarrow
\infty$ we tend to loose both these properties since we will be obliged to
take $a_T \rightarrow + \infty$. In order to obtain uniform upper and lower
bounds we proceed as follows.

Denote by $\mathbbm{H}_a$ the space of progressively measurable processes
which are $\mathbbm{P}$--almost surely in $L^2 (\mathbbm{R}_+ \times
\Lambda)$. We say that an element $v$ of $\mathbbm{H}_a$ is a
\tmtextit{drift}. Consider the measure $\mathbbm{Q}_T$ on $\left( \Omega, \CB
\right)$ whose Radon--Nykodim derivative wrt. $\mathbbm{P}$ is given by
\[ \frac{\mathd \mathbbm{Q}_T}{\mathd \mathbbm{P}} = \frac{e^{- V_T
   (Y_T)}}{\CZ_T} . \]
Since $Y_T$ depends on finitely many Brownian motions $(B^n)_n$ then it is
well known~{\cite{revuz_continuous_2004,follmer_entropy_1985}} that any
$\mathbbm{P}$--absolutely continuous probability can be expressed via Girsanov
transform. In particular, by the Brownian martingale representation theorem
there exists a drift $u^T \in \mathbbm{H}_a$ such that
\[ \frac{\mathd \mathbbm{Q}_T}{\mathd \mathbbm{P}} = \exp \left(
   \int_0^{\infty} u_s^T \mathd X_s - \frac{| \Lambda |}{2} \int_0^{\infty} \|
   u_s^T \|^2 \mathd s \right), \]
and the entropy of $\mathbbm{Q}_T$ wrt. $\mathbbm{P}$ is given by
\[ H (\mathbbm{Q}_T |\mathbbm{P}) =\mathbbm{E}_{\mathbbm{Q}_T} \left[ \log
   \frac{\mathd \mathbbm{Q}_T}{\mathd \mathbbm{P}} \right] = \frac{| \Lambda
   |}{2} \mathbbm{E}_{\mathbbm{Q}_T} \left[ \int_0^{\infty} \| u_s^T \|^2
   \mathd s \right] . \]
Here equality holds also if one of the two quantities is $+ \infty$. By
Girsanov theorem, the canonical process $X$ is a semimartingale under
$\mathbbm{Q}_T$ with decomposition
\[ X_t = \tilde{X}_t + \int_0^t u^T_s \mathd s, \]
where $(\tilde{X}_t)_t$ is a cylindrical $\mathbbm{Q}_T$--Brownian motion in
$L^2 (\Lambda)$. This implies that, under $\mathbbm{Q}_T$ the process
$(Y_t)_t$ has the semimartingale decomposition $Y_t = W_t + U_t$ with
\[ W_t \assign \int_0^t J_s \mathd X_s, \quad \text{and} \quad U_t = I_t
   (u^T), \]
where for any drift $v \in \mathbbm{H}_a$ we define $I_t (v) \assign \int_0^t
J_s v_s \mathd s$. The integral in the density can be restricted to $[0, T]$
since $u_t^T = 0$ if $t > T$. Now
\begin{equation}
  - \log \CZ_T = - \log \left[ e^{- V_T (Y_T)} \left( \frac{\mathd
  \mathbbm{Q}_T}{\mathd \mathbbm{P}} \right)^{- 1} \right] = V_T (Y_T) +
  \int_0^{\infty} u_s^T \mathd X_s - \frac{| \Lambda |}{2} \int_0^{\infty} \|
  u_s^T \|^2 \mathd s, \label{eq:equilib}
\end{equation}
and taking expectation of~(\ref{eq:equilib}) wrt $\mathbbm{Q}_T$ we get
\begin{equation}
  - \log \CZ_T =\mathbbm{E}_{\mathbbm{Q}_T} \left[ V_T (W_T + I_T (u^T)) +
  \frac{| \Lambda |}{2} \int_0^{\infty} \| u_s^T \|^2 \mathd s \right] .
  \label{eq:pressure}
\end{equation}
This will be our key formula to get an upper bound on $Z_T$. For any $v \in
\mathbbm{H}_a$ define the measure $\mathbbm{Q}^v$ by
\[ \frac{\mathd \mathbbm{Q}^v}{\mathd \mathbbm{P}} = \exp \left(
   \int_0^{\infty} v_s \mathd X_s - \frac{| \Lambda |}{2} \int_0^{\infty} \|
   v_s \|^2 \mathd s \right) . \]
Denote with $\mathbbm{H}_c \subseteq \mathbbm{H}_a$ the set of drifts for
which $\mathbbm{Q}^v (\Omega) = 1$. In particular $u^T \in \mathbbm{H}_c$. By
Jensen's inequality and Girsanov transformation we have
\[ - \log \CZ_T = - \log \mathbbm{E}_{\mathbbm{P}} [e^{- V_T (Y_T)}] = - \log
   \mathbbm{E}^v \left[ e^{- V_T (Y_T) - \int_0^{\infty} v_s \mathd X_s +
   \frac{| \Lambda |}{2} \int_0^{\infty} \| v_s \|^2 \mathd s} \right] \]
\[ \leqslant \mathbbm{E}^v \left[ V_T (Y_T) + \int_0^{\infty} v_s \mathd X_s -
   \frac{| \Lambda |}{2} \int_0^{\infty} \| v_s \|^2 \mathd s \right] . \]
for all $v \in \mathbbm{H}_c$, where $\mathbbm{E}^v \assign
\mathbbm{E}_{\mathbbm{Q}^v}$. We conclude that
\begin{equation}
  - \log \CZ_T \leqslant \mathbbm{E}^v \left[ V_T (W^v_T + I_T (v)) + \frac{|
  \Lambda |}{2} \int_0^{\infty} \| v_s \|^2 \mathd s \right]
  \label{eq:pressure-bound}
\end{equation}
where $Y = W^v_T + I_T (v)$ and $\tmop{Law}_{\mathbbm{Q}^v} (W^v) =
\tmop{Law}_{\mathbbm{P}} (Y)$. This bound is saturated when $v = u^T$. We
record this result in the following lemma which is a precursor of our main
tool to obtain bounds on the partition function and related objects.

\begin{lemma}
  \label{lemma:variational-easy}The following variational formula for the
  pressure $- \log \CZ_T$ holds:
  \[ - \frac{1}{| \Lambda |} \log \CZ_T = - \frac{1}{| \Lambda |} \log
     \mathbbm{E} [e^{- V_T (Y_T)}] = \min_{v \in \mathbbm{H}_c} \mathbbm{E}^v
     \left[ \frac{1}{| \Lambda |} V_T (W^v_T + I_T (v)) + \frac{1}{2}
     \int_0^{\infty} \| v_s \|^2_{L^2} \mathd s \right] . \]
\end{lemma}

This formula is nice and easy to prove but somewhat incovenient for certain
manipulations since the space $\mathbbm{H}_c$ is indirectly defined and the
reference measure $\mathbbm{E}^v$ depends on the drift $v$. A more
straighforward formula has been found by
Bou{\'e}--Dupuis~{\cite{boue_variational_1998}} which involves a the fixed
canonical measure $\mathbbm{P}$ and a general adapted drift $u \in
\mathbbm{H}_a$. This formula will be our main tool in the following.

\begin{theorem}
  \label{th:variational}The Bou{\'e}--Dupuis variational formula for the
  pressure $- \log \CZ_T$ holds:
  \[ - \frac{1}{| \Lambda |} \log \CZ_T = - \frac{1}{| \Lambda |} \log
     \mathbbm{E} [e^{- V_T (Y_T)}] = \inf_{v \in \mathbbm{H}_a} \mathbbm{E}
     \left[ \frac{1}{| \Lambda |} V_T (Y_T + I_T (v)) + \frac{| \Lambda |}{2}
     \int_0^{\infty} \| v_s \|^2_{L^2} \mathd s \right] . \]
  where the expectation is taken wrt to the measure $\mathbbm{P}$ on $\Omega$.
\end{theorem}

\begin{proof}
  The original proof can be found in
  Bou{\'e}--Dupuis~{\cite{boue_variational_1998}} for functionals bounded
  above. In our setting the formula can be proved using the result of
  {\"U}st{\"u}nel~{\cite{ustunel_variational_2014}} by observing that $V_T
  (Y_T)$ is a \tmtextit{tame} functional, according to his definitions.
  Namely, for some $p, q \geqslant 1$ such that $1 / p + 1 / q = 1$ we have
  \[ \mathbbm{E} [| V_T (Y_T) |^p] +\mathbbm{E} [e^{- q V_T (Y_T)}] < + \infty
     . \]
  
\end{proof}

\begin{remark}
  Some observations on this variational formulas.
  \begin{enumeratealpha}
    \item They originates directly from the variational formula for the free
    energy of a statistical mechanics systems: $V_T$ playing the role of the
    internal energy and the quadratic term playing the role of the entropy.
    
    \item Note that the infimum could not be attained in
    Theorem~\ref{th:variational} while it is attained in
    Lemma~\ref{lemma:variational-easy}.
    
    \item The drift generated by absolutely continuous perturbations of the
    Wiener measure has been introduced and studied by
    F{\"o}llmer~{\cite{follmer_entropy_1985}}.
    
    \item The Bou{\'e}--Dupuis formula in Theorem~\ref{th:variational} is
    easier to use than the formula in Lemma~\ref{lemma:variational-easy} since
    the probability do not depend on the drift $v$. However going from one
    formulation to the other requires proving that certain SDEs with
    functional drift admits strong solutions and that one is able to
    approximate unbounded functionals $V_T$ by bounded ones. This requires the
    notion of \tmtextit{tame} functional, introduce by
    {\"U}st{\"u}nel~{\cite{ustunel_variational_2014}}.
    
    \item The Bou{\'e}--Dupuis formula is a non--Markovian and infinite
    dimensional extension of the well known stochastic control problem
    representation of the Hamilton--Jacobi--Bellman equation in finite
    dimensions~{\cite{fleming_controlled_2005}}.
    
    \item See {\"U}st{\"u}nel~{\cite{ustunel_variational_2014}} and
    Lehec~{\cite{lehec_representation_2013}} for a streamlined proof of the
    Bou{\'e}--Dupuis formula and for applications of the formula to functional
    inequalities on Gaussian measures. For example, from this formula is not
    difficult to prove integrability of functionals which are Lipshitz in the
    Cameron--Martin directions.
  \end{enumeratealpha}
\end{remark}

The next lemma provides a deterministic regularity result for $I (v)$ which
will be useful below. It says that the drift $v$ generates shifts of the
Gaussian free field in directions which belong to $H^1$ uniformly in the scale
parameter. The space $H^1$ is the Cameron--Martin space of the free
field~{\cite{janson_gaussian_1997}}.

\begin{lemma}
  \label{lemma:reg-drift}For any $v \in L^2 (\mathbbm{R}_+ \times \Lambda)$ we
  have
  \[ \sup_{0 \leqslant t \leqslant T} \| I_t (v) \|_{H^1}^2 \lesssim \int_0^T
     \| v_s \|^2_{L^2} \mathd s. \]
\end{lemma}

\begin{proof}
  Using the fact that $\sigma_s (\mathD)$ is diagonal in Fourier space, and
  denoting with $(e_k)_{k \in \mathbbm{Z}^d}$ the basis of trigonometric
  polynomials, we have
  \[ \langle \int_0^t \sigma_s (\mathD) v_s \mathd s, \int_0^t \sigma_s
     (\mathD) v_s \mathd s \rangle = \sum_k \left| \int_0^t \langle \sigma_s
     (\mathD) e_k, v_s \rangle \mathd s \right|^2 \]
  \[ \leqslant \sum_k \left( \int_0^t | \langle \sigma_s (\mathD) e_k, e_k
     \rangle |^2 \mathd s \right) \left( \int_0^t | \langle e_k, v_s \rangle
     |^2 \mathd s \right) \leqslant | \Lambda | \int_0^t \| v_s \|^2_{L^2}
     \mathd s \sup_k \int_0^t \langle e_k, \sigma_s (\mathD)^2 e_k \rangle
     \mathd s \]
  \[ \leqslant | \Lambda | \int_0^t \| v_s \|^2_{L^2} \mathd s \sup_k \langle
     e_k, \rho_t^2 (\mathD) e_k \rangle \leqslant | \Lambda | \int_0^t \| v_s
     \|^2_{L^2} \mathd s. \]
  Therefore
  \[ \| I_t (v) \|_{H^1}^2 \lesssim \| \langle \mathD \rangle I_t (v)
     \|_{L^2}^2 = \frac{1}{| \Lambda |} \langle \int_0^t \sigma_s (\mathD) v_s
     \mathd s, \int_0^t \sigma_s (\mathD) v_s \mathd s \rangle \leqslant
     \int_0^t \| v_s \|^2_{L^2} \mathd s. \]
\end{proof}

\begin{notation}
  In the estimates below the symbol $E (\lambda)$ will denote a generic
  positive deterministic quantity, not depending on $| \Lambda |$ and such
  that $E (\lambda) / \lambda^3 \rightarrow 0$ as $\lambda \rightarrow 0$.
  Moreover the symbol $Q_T$ will denote a generic random variable measurable
  wrt. $\sigma ((W_t)_{t \in [0, T]})$ and belonging to $L^1 (\mathbbm{P})$
  uniformly in $T$. 
\end{notation}

\section{Two dimensions}\label{sec:two-d}

As a warm up we consider here the case $d = 2$. From
Theorem~\ref{th:variational} we see that the relevant quantity to bound is of
the form
\begin{equation}
  \Theta_T (u) \assign \mathbbm{E} \left[ \frac{1}{| \Lambda |} V_T (W_T + I_T
  (u)) + \frac{1}{2} \int_0^{\infty} \| u_s \|_{L^2 (\Lambda, \mathd x)}^2
  \mathd s \right], \label{eq:var-energy}
\end{equation}
for $u \in \mathbbm{H}_a$. Let $Z_t = I_t (u)$ and let $W_t = Y_t$ as a
mnemonic of the fact that under $\mathbbm{P}$ the process $W$ is a martingale.
From now on we leave implicit the integration variable over the spatial domain
$\Lambda$. Choosing
\begin{equation}
  a_T = 6\mathbbm{E} [W_T (\xi)^2], \qquad b_T = 3\mathbbm{E} [W_T (\xi)^2]^2,
  \label{eq:renorm-const-2d}
\end{equation}
we have
\[ V_T (W_T + Z_T) = \int_{\Lambda} \lambda \llbracket W_T^4 \rrbracket +
   \int_{\Lambda} 4 \lambda \llbracket W_T^3 \rrbracket Z_T + \int_{\Lambda} 6
   \lambda \llbracket W_T^2 \rrbracket Z_T^2 + \int_{\Lambda} 6 \lambda W_T
   Z_T^3 + \int_{\Lambda} \lambda Z_T^4, \]
where
\[ \llbracket W_T^4 \rrbracket (\xi) \assign W_T (\xi)^4 - 6\mathbbm{E} [W_T
   (\xi)^2] W_T (\xi)^2 + 3\mathbbm{E} [W_T (\xi)^2]^2, \]
\[ \llbracket W_T^3 \rrbracket (\xi) \assign W_T (\xi)^3 - 3\mathbbm{E} [W_T
   (\xi)^2] W_T (\xi), \]
\[ \llbracket W_T^2 \rrbracket (\xi) \assign W_T (\xi)^2 -\mathbbm{E} [W_T
   (\xi)^2], \]
denote the Wick powers of the Gaussian r.v. $W_T
(\xi)$~{\cite{janson_gaussian_1997}}. The reader can check the well known
property of these polynomials when seen as stochastic processes in $T$, namely
that they are $\mathbbm{P}$--martingales wrt. the filtration of $(W_t)_t$. In
particular they have an expression as iterated stochastic integrals wrt. the
Brownian motions $(B^n_t)_{t, n}$ introduced in eq.~(\ref{eq:rep-Y}). Using
Theorem~\ref{th:variational} with $u = 0$ we readily have an upper bound for
the pressure:
\[ - \frac{1}{| \Lambda |} \log \CZ_T \leqslant \lambda \mathbbm{E} \left[
   \bint \llbracket W_T^4 \rrbracket \right] = 0. \]
As for an effective lower bound we need to estimate from below the average
under $\mathbbm{P}$ of the variational expression
\[ \lambda \bint \llbracket W_T^4 \rrbracket + 4 \lambda \bint \llbracket
   W_T^3 \rrbracket Z_T + 6 \lambda \bint \llbracket W_T^2 \rrbracket Z_T^2 +
   4 \lambda \bint W_T Z_T^3 + \lambda \bint  Z_T^4 + \frac{1}{2} \int_0^T
   \| u_s \|^2_{L^2} \mathd s. \]
The strategy we adopt is to bound \tmtextit{pathwise}, and for a generic drift
$u$, the contributions
\[ \Psi_T (u) \assign \underbrace{4 \lambda \bint \llbracket W_T^3 \rrbracket
   Z_T}_{\Iota} + \underbrace{6 \lambda \bint \llbracket W_T^2 \rrbracket
   Z_T^2}_{\Iota \Iota} + \underbrace{4 \lambda \bint W_T Z_T^3}_{\Iota \Iota
   \Iota}, \]
in term of quantities involving only the Wick powers of $W$ which we can
control in expectation and the last two \tmtextit{positive terms}
\[ \Phi_T (u) \assign \frac{1}{2} \int_0^T \| u_s \|^2_{L^2} \mathd t +
   \lambda \bint Z_T^4 . \]
Any residual positive contribution depending on $u$ can be dropped in the
lower bound making the dependence on the drift disappear. To control term
$\Iota$ we see that by duality and Young's inequality, for any $\delta > 0$,
\begin{equation}
  \left| 4 \lambda \bint \llbracket W_T^3 \rrbracket Z_T \right| \leqslant 4
  \lambda \| \llbracket W_T^3 \rrbracket \|_{H^{- 1}} \| Z_T \|_{H^1}
  \leqslant C (\delta, d) \lambda^2 \| \llbracket W_T^3 \rrbracket \|_{H^{-
  1}}^2 + \delta \int_0^T \| u_s \|^2_{L^2} \mathd s. \label{eq:term-i}
\end{equation}
For the term $\Iota \Iota$ the following fractional Leibniz rule is of help:

\begin{proposition}
  \label{FractionalLeibniz}Let $1 < p < \infty$ and $p_1, p_2 > 1$ such that
  $\frac{1}{p_1} + \frac{1}{p_2} = \frac{1}{p}$. Then for every such $p_, p_1,
  p_2$ and every $s, \alpha \geq 0$ there exists a constant C such that
  \[ \| \langle \mathD \rangle^s (f g) \|_{L^p} \leq \| \langle \mathD
     \rangle^{s + \alpha} f \|_{L^{p_2}}  \| \langle \mathD \rangle^{- \alpha}
     g \|_{L^{p_1}} + \| \langle \mathD \rangle^{s + \alpha} g \|_{L^{p_1}} 
     \| \langle \mathD \rangle^{- \alpha} f \|_{L^{p_2}} . \]
\end{proposition}

\begin{proof}
  See~{\cite{gulisashvili_exact_1996}}.
\end{proof}

Using Proposition~\ref{FractionalLeibniz} we get, for any $\delta > 0$,
\begin{equation}
  \begin{array}{lll}
    \left| 6 \lambda \bint \llbracket W_T^2 \rrbracket Z_T^2  \right| &
    \lesssim & \lambda \| \llbracket W_T^2 \rrbracket \|_{W^{- \varepsilon,
    5}} \| Z_T^2 \|_{W^{\varepsilon, \frac{5}{4}}}\\
    & \lesssim & \lambda \| \llbracket W_T^2 \rrbracket \|_{W^{- \varepsilon,
    5}} \| Z_T \|_{W^{\varepsilon, 2}} \| Z_T \|_{L^{\frac{10}{3}}}\\
    & \lesssim & \lambda \| \llbracket W_T^2 \rrbracket \|_{W^{- \varepsilon,
    5}} \| Z_T \|_{W^{1, 2}} \| Z_T \|_{L^4}\\
    & \leqslant & \frac{C^2 \lambda^3}{2 \delta} \| \llbracket W_T^2
    \rrbracket \|^4_{W^{- \varepsilon, 5}} + \frac{\delta}{4} \| Z_T
    \|^2_{W^{1, 2}} + \frac{\delta \lambda}{4} \| Z_T \|_{L^4}^4 .
  \end{array} \label{eq:term-ii}
\end{equation}


In order to bound the term $\Iota \Iota \Iota$ we observe the following:

\begin{lemma}
  \label{lemma:bound-cubic}Let $f \in W^{- 1 / 2 - \varepsilon, p}$ for every
  $1 \leqslant p < \infty$ and $\varepsilon > 0$. We claim that for any
  $\delta > 0$ there exists a constant $C = C (\delta, d)$, and an exponent $K
  < \infty$ such that for any $g \in W^{1, 2}$
  \[ \lambda \left| \bint f g^3 \right| \leqslant E (\lambda) \| f \|^K_{W^{-
     1 / 2 - \varepsilon, p} } + \delta (\| g \|^2_{W^{1 - \varepsilon, 2}} +
     \lambda \| g \|^4_{L^4 }) . \]
\end{lemma}

\begin{proof}
  By duality $\left| \bint f g^3 \right| \leqslant \| f \|_{W^{- 1 / 2 -
  \varepsilon, p}} \| g^3 \|_{W^{1 / 2 + \varepsilon, p'}}$ . Applying again
  Proposition \ref{FractionalLeibniz}, we get
  \begin{eqnarray*}
    \| g^3 \|_{W^{1 / 2 + \varepsilon, p'}} & \lesssim & \| \langle \mathD
    \rangle^{1 / 2 + \delta} g^3 \|_{L^{p'}} \lesssim \| \langle \mathD
    \rangle^{1 / 2 + \delta} g\|_{L^{p_1}} \| g \|_{L^4}^2\\
    & \lesssim & \| \langle \mathD \rangle^{1 / 2 + 2 \delta} g\|_{L^2
    } \| g \|_{L^4}^2 \lesssim \|g\|^{\frac{1 / 2 + 2 \delta}{1 -
    \delta}}_{H^{1 - \delta}} \| g \|_{L^2}^{\frac{1 / 2 - 3 \delta}{1 -
    \delta}} \| g \|_{L^4}^2\\
    & \lesssim & \|g\|^{\frac{1 / 2 + 2 \delta}{1 - \delta}}_{H^{1 -
    \delta}} \| g \|_{L^4}^{\frac{5 / 2 - 5 \delta}{1 - \delta}} .
  \end{eqnarray*}
  So
  \begin{eqnarray*}
    \lambda \left| \bint f g^3 \right| & \lesssim & \lambda^{\frac{3 / 2 +
    \delta}{4 (1 - \delta)}} \| f \|_{W^{- 1 / 2 - \varepsilon, p}}
    \|g\|^{\frac{1 / 2 + 2 \delta}{1 - \delta}}_{H^{1 - \delta}}
    \lambda^{\frac{5 / 2 - 5 \delta}{4 (1 - \delta)}} \| g
    \|_{L^4}^{\frac{5 / 2 - \delta}{1 - \delta}}\\
    & \lesssim & \lambda^{\frac{3 / 2 + \delta}{(1 - \delta) (1 / 2 -
    \delta)}} \| f \|^{\frac{4}{1 / 2 - \delta}}_{W^{- 1 / 2 - \varepsilon,
    p}} + \delta (\|g\|^2_{H^{1 - \delta}} + \lambda \| g \|_{L^4}^4) .
  \end{eqnarray*}
  If $\delta$ is chosen small enough, $p$ is sufficiently large and $p_1$ is
  chosen suffiently close to $2$.
\end{proof}

\begin{remark}
  For the $d = 2$ case it would have been enough to estimate $\| g^3
  \|_{W^{\varepsilon, p}}$. The stronger estimate will be useful for $d = 3$
  since there we will only have $W_T \in W^{- 1 / 2 - \varepsilon, p}$.
\end{remark}

Then
\begin{equation}
  \left| 4 \lambda \bint W_T Z_T^3 \right| \leqslant C E (\lambda) \| W_T
  \|^K_{W^{- 1 / 2 - \varepsilon, p} } + \delta (\| Z_T \|^2_{W^{1 -
  \varepsilon, 2}} + \lambda \| Z_T \|^4_{L^4 }), \label{eq:term-iii}
\end{equation}
and using eqs.~(\ref{eq:term-i}),~(\ref{eq:term-ii}) and~(\ref{eq:term-iii})
together we obtain for $\delta$ small enough
\begin{equation}
  | \Psi_T (u) | \leqslant O (\lambda^2) [1 + \| \llbracket W_T^3 \rrbracket
  \|_{H^{- 1}}^2 + \| \llbracket W_T^2 \rrbracket \|^4_{W^{- \varepsilon, 5}}
  + \| W_T \|^K_{W^{- 1 / 2 - \varepsilon, p} }] + \delta \Phi_T (u) .
  \label{eq:bound-Psi}
\end{equation}
and therefore
\[ \Theta_T (u) \geqslant - C\mathbbm{E} [1 + \| \llbracket W_T^3 \rrbracket
   \|_{H^{- 1}}^2 + \| \llbracket W_T^2 \rrbracket \|^4_{W^{- \varepsilon, 5}}
   + \| W_T \|^K_{W^{- 1 / 2 - \varepsilon, p} }] . \]
Using eq.~(\ref{eq:term-i}), eq.~(\ref{eq:term-ii}) and
Lemma~\ref{lemma:bound-cubic} in eq.~(\ref{eq:var-energy}) we get
\begin{eqnarray*}
  \Theta_T (u) & \geqslant & - C\mathbbm{E} [1 + \| \llbracket W_T^3
  \rrbracket \|_{H^{- 1}}^2 + \| \llbracket W_T^2 \rrbracket \|^4_{W^{-
  \varepsilon, 5}} + \| W_T \|^K_{W^{- 1 / 2 - \varepsilon, p} }] +
  \frac{1}{2} \mathbbm{E} [\Phi_T (u)]\\
  & \geqslant & - C\mathbbm{E} [1 + \| \llbracket W_T^3 \rrbracket \|_{H^{-
  1}}^2 + \| \llbracket W_T^2 \rrbracket \|^4_{W^{- \varepsilon, 5}} + \| W_T
  \|^K_{W^{- 1 / 2 - \varepsilon, p} }] .
\end{eqnarray*}
This last average do not depends anymore on the drift and we are only left to
show that
\[ \sup_T \mathbbm{E} [\| \llbracket W_T^3 \rrbracket \|^2_{H^{- 1}} + \|
   \llbracket W_T^2 \rrbracket \|^4_{W^{- \varepsilon, 5}} + \| W_T \|^K_{W^{-
   1 / 2 - \varepsilon, p}}] < \infty . \]
However it is well known that the Wick powers of the two dimensional Gaussian
free field are distributions belonging to $L^a (\Omega, W^{- \varepsilon, b})$
for any $a \geqslant 1$ and $b \geqslant 1$ and hypercontractivity plus an
easy argument gives the uniform boundedness of the above averages. Therefore
we have established that

\begin{theorem}
  \label{th:twod-bound}For any $\lambda > 0$ we have
  \[ \sup_T \frac{1}{| \Lambda |} \left| \log \CZ_T \right| \lesssim O
     (\lambda^2), \]
  where the constant in the r.h.s. is independent of $\Lambda$. 
\end{theorem}

\begin{remark}
  Note this also implies that $e^{- V_T (Y_T)}$ is in all the $L^p$ spaces
  wrt. the measure $\mathbbm{P}$ uniformly in $T$ and for any $p \geqslant 1$
  since the above argument remains valid upon replacing $\lambda$ with
  $\lambda p$.
\end{remark}

\begin{remark}
  With slightly different ways to estimate the various terms it would be
  possible to show that the bound in Theorem~\ref{th:twod-bound} is actually
  $o (\lambda^3)$ as in the $d = 3$ case. 
\end{remark}

\section{Three dimensions}\label{sec:three-d}

We would like now to pursue a similar strategy when $d = 3$ where the
regularity of the Wick products is less pleasant. In particular we only have
$\llbracket W_T^3 \rrbracket \in \VV^{- 3 / 2 - \kappa}$, $\llbracket W_T^2
\rrbracket \in \VV^{- 1 - \kappa}$ uniformly in $T$ for any small $\kappa >
0$. This implies that there is no straighforward way to control the terms
\begin{equation}
  \bint \llbracket W_T^3 \rrbracket Z_T, \qquad \text{and} \qquad \bint
  \llbracket W_T^2 \rrbracket Z_T^2, \label{eq:mixed-terms}
\end{equation}
like we did in Section~\ref{sec:two-d}. Indeed the only apriori estimate on
the regularity of $Z_T = I_T (u)$ is in $H^1$, coming from
Lemma~\ref{lemma:reg-drift} and the quadratic term in the variational
functional $\Theta_T (u)$. It is also well known that in three dimensions
there are further divergences beyond the Wick ordering which have to be
subtracted in order for the limiting measure to be non-trivial. For these
reasons we introduce further scale dependent renormalization constants
$\gamma_T, \delta_T$ into the energy $V_T$ beyond Wick ordering and let
\[ \frac{1}{| \Lambda |} V_T (Y_T) = \bint (\lambda \llbracket Y_T^4
   \rrbracket - \gamma_T \llbracket Y_T^2 \rrbracket - \delta_T) . \]
Repeating the computation from Section~\ref{sec:two-d} we arrive at
\begin{eqnarray*}
  \frac{\Theta_T (u)}{| \Lambda |} & = & \mathbbm{E} \left[ \lambda \bint
  \mathbbm{W}_T^3 Z_T + \frac{\lambda}{2} \bint \mathbbm{W}^2_T Z_T^2 + 4
  \lambda \bint W_T Z_T^3 \right]\\
  &  & -\mathbbm{E} \left[ 2 \gamma_T \bint W_T Z_T + \gamma_T \bint
  Z_T^2 \right] +\mathbbm{E}^u \left[ \lambda \bint Z_T^4 + \frac{1}{2}
  \int_0^T \| u_s \|^2_{L^2} \mathd s \right],
\end{eqnarray*}
where we introduced the convenient notations
\[ \mathbbm{W}_t^3 \assign 4 \llbracket W_t^3 \rrbracket, \qquad
   \mathbbm{W}^2_t \assign 12 \llbracket W_t^2 \rrbracket . \]
The rest of this section will be concerned with lower and upper bounds of this
quantity leading to the proof of Theorem~\ref{th:main-3d}.

\subsection{Lower bound}

\subsubsection{Bound for the linear term}

We are going to use the quadratic term $\int_0^T \| u_s \|^2_{L^2} \mathd s$
to absorb the mixed terms~(\ref{eq:mixed-terms}). In order to do so we need to
represent them as quantities flowing along the scale parameter $t$. This will
allow to identify suitable counterterms. By Ito's formula
\[ \lambda \bint \mathbbm{W}_T^3 Z_T = \lambda \int_0^T \bint
   \mathbbm{W}_t^3 \dot{Z}_t \mathd t + \text{martingale}, \]
where we used that $(\mathbbm{W}_t^3)_t$ is a martingale and that $(Z_t)_t$ is
of bounded variation. We now introduce a new drift $w$ via the formula
\begin{equation}
  u_t = - \lambda J_t \mathbbm{W}_t^3 + w_t, \label{eq:ansatz-1}
\end{equation}
we have
\begin{equation}
  \label{changeeq} Z_s = - \lambda \int_0^s J^2_t \mathbbm{W}_t^3 \mathd t
  + \int_0^s J_t w_t \mathd t = - \lambda \mathbbm{W}_s^{[3]} + K_s,
\end{equation}
where
\[ K_t \assign I_t (w), \qquad \text{and} \qquad \mathbbm{W}_t^{[3]} \assign
   \int_0^t J^2_s \mathbbm{W}_s^3 \mathd s. \]
Note that $\dot{Z_t} = - \lambda J^2_t \mathbbm{W}_t^3 + J_t w_t$. It holds
\[ \lambda \bint \mathbbm{W}_T^3 Z_T + \frac{1}{2} \int_0^T \| u_s \|^2_{L^2}
   \mathd s = - \frac{1}{2} \lambda^2 \int_0^T \bint (J_t \mathbbm{W}_t^3)^2
   \mathd t + \frac{1}{2} \int_0^T \| w_s \|^2_{L^2} \mathd s +
   \tmop{martingale} . \]
Here we have eliminated the mixed divergent term $\bint \mathbbm{W}_T^3 Z_T$
and replaced it with a divergent but purely stochastic term $\int_0^T \bint
(J_t \mathbbm{W}_t^3)^2 \mathd t$ which does not affect anymore the
variational problem and can be explicitly removed by adding its average to
$\delta_T$. This implies in turn that we are not able to control $(Z_t)_t$ in
$H^1$ anymore and we should rely on the relation~(\ref{changeeq}) and on a
control over the $H^1$ norm of $(K_t)_t$ coming from the residual quadratic
term $\int_0^T \| w_s \|^2_{L^2} \mathd s$.

\begin{remark}
  This lack of uniform regularity of the full drift $u$ is a sign of the fact
  that the limiting measure $\nu$ is expected not to be absolutely continuous
  wrt. to the Gaussian free field. We do not make this connection rigorous in
  this paper and leave the discussion to a further publication.
\end{remark}

\subsubsection{Bound for the quadratic term}

Taking into account the relation~(\ref{changeeq}) we write
\[ \frac{\lambda}{2} \bint \mathbbm{W}^2_T Z_T^2 = \frac{\lambda^3}{2} \bint
   \mathbbm{W}^2_T (\mathbbm{W}_T^{[3]})^2 - \lambda^2 \bint \mathbbm{W}^2_T
   \mathbbm{W}_T^{[3]} K_T + \frac{\lambda}{2} \bint \mathbbm{W}^2_T K_T^2 .
\]
The first term is already purely stochastic and will give a contribution
independent of $u$ which can again be absorbed in $\delta_T$. \ In order to
proceed we will isolate the divergent parts of the product via paracontrolled
calculus (see Appendix~\ref{sec:appendix-para} for details). This is the
content of the following

\begin{lemma}
  For any small $\varepsilon > 0$ there exists $\delta > 0$ such that
  \begin{eqnarray*}
    &  & \left| - \lambda^2 \bint \mathbbm{W}^2_T \mathbbm{W}_T^{[3]} K_T +
    \frac{\lambda}{2} \bint \mathbbm{W}^2_T K_T^2 - \lambda \bint
    (\mathbbm{W}^2_T \succ Z_T) K_T + \lambda^2 \bint (\mathbbm{W}^2_T \circ
    \mathbbm{W}_T^{[3]}) K_T \right|\\
    &  & \leqslant C (\varepsilon, \delta) E (\lambda) Q_T + \varepsilon \|
    K_T \|^2_{H^{1 - \delta}} + \varepsilon \lambda \| Z_T \|^4_{L^4} .
  \end{eqnarray*}
\end{lemma}

\begin{proof}
  We compute
  \begin{eqnarray*}
    - \lambda^2 \bint \mathbbm{W}^2_T \mathbbm{W}_T^{[3]} K_T +
    \frac{\lambda}{2} \bint \mathbbm{W}^2_T K_T^2 & = & \lambda \bint
    (\mathbbm{W}^2_T \succ Z_T) K_T - \lambda^2 \bint (\mathbbm{W}^2_T \circ
    \mathbbm{W}_T^{[3]}) K_T\\
    &  & - \lambda^2 \bint (\mathbbm{W}^2_T \prec \mathbbm{W}_T^{[3]}) K_T +
    \frac{\lambda}{2} \bint (\mathbbm{W}^2_T \prec K_T) K_T\\
    &  & + \frac{\lambda}{2} \left( \bint (\mathbbm{W}^2_T \circ K_T) K_T -
    \bint (\mathbbm{W}^2_T \succ K_T) K_T \right) .
  \end{eqnarray*}
  By Proposition~\ref{adjointparaproduct},
  \begin{eqnarray*}
    \lambda \left| \bint (\mathbbm{W}^2_T \succ K_T) K_T  -  \bint
    (\mathbbm{W}^2_T \circ K_T) K_T \right| &  \lesssim &  \lambda \|
    \mathbbm{W}^2_T \|_{\VV^{- 1 - \delta}}  \| K_T \|^2_{H^{1 / 2 +
    \delta}}\\
    & \lesssim & \lambda^{\frac{3 / 2 - \delta}{2 (1 - \delta)}} \|
    \mathbbm{W}^2_T \|_{\VV^{- 1 - \delta}} \times \\ & & \times \| K_T \|^{2 \left( \frac{1 / 2 +
    \delta}{1 - \delta} \right)}_{H^{1 - \delta}} \lambda^{\frac{1 / 2 -
    \delta}{2 (1 - \delta)}} \| K_T \|^{2 \left( \frac{1 / 2 - \delta}{1 -
    \delta} \right)}_{L^2}\\
    & \leqslant & C \lambda^{^{\frac{3 / 2 - \delta}{1 / 2 - 3 \delta}}} \|
    \mathbbm{W}^2_T \|^{\left( 1 - \left( \frac{3 / 2 + \delta}{2 (1 -
    \delta)} \right) \right)^{- 1}}_{\VV^{- 1 - \delta}} + \varepsilon \| K_T
    \|^2_{H^{1 - \delta}}\\
    &  & + \varepsilon \lambda \| K_T \|^4_{L^2} .
  \end{eqnarray*}
  Furthermore by Proposition~\ref{paraproductestimate},
  \begin{eqnarray*}
    \left| \lambda \bint (\mathbbm{W}^2_T \prec K_T) K_T \right| & \lesssim &
    \lambda \| \mathbbm{W}^2_T \|_{\VV^{- 1 - \delta}}  \| K_T \|^2_{H^{1 /
    2 + \delta}}\\
    & \leqslant & C \lambda^{\frac{3 / 2 - \delta}{1 / 2 + \delta}} \|
    \mathbbm{W}^2_T \|^{\left( 1 - \left( \frac{3 / 2 + \delta}{2 (1 -
    \delta)} \right) \right)^{- 1}}_{\VV^{- 1 - \delta}} + \varepsilon \| K_T
    \|^2_{H^{1 - \delta}} + \varepsilon \lambda \| K_T \|^4_{L^2},
  \end{eqnarray*}
  and
  \begin{eqnarray*}
    \left| \lambda^2 \bint (\mathbbm{W}^2_T \prec \mathbbm{W}_T^{[3]}) K_T
    \right| & \lesssim & \lambda^2 \| \mathbbm{W}^2_T \|_{\VV^{- 1 - \delta /
    2}}  \| \mathbbm{W}_T^{[3]} \|_{\VV^{^{1 / 2 - \delta / 2}}} \| K_T
    \|_{H^{1 / 2 + \delta}}\\
    & \leqslant & C (\delta) \lambda^4 \left( \| \mathbbm{W}^2_T \|_{\VV^{- 1
    - \delta / 2}}  \| \mathbbm{W}_T^{[3]} \|_{\VV^{^{1 / 2 - \delta / 2}}}
    \right)^2 + \delta \| K_T \|^2_{H^{1 / 2 + \delta}} .
  \end{eqnarray*}
\end{proof}

The resonant term $\mathbbm{W}^2_T \circ \mathbbm{W}_T^{[3]}$ will behave well
enough after suitable renormalization, as we will see later. As far as the
term $\lambda \bint (\mathbbm{W}^2_T \succ Z_T) K_T$ in concerned, we want to
absorb this into $\int \| w_s \|^2 \mathd s$ like we did with the linear term.
Before we can do this we must make sure that after applying Ito's formula it
still possible to use $\bint Z^4_T$ to control some of the growth of this term
since we already have quadratic (critical) dependence on $(K_t)_t$ which
cannot be fully taken care of by $\int \| w_s \|^2 \mathd s$. Choose $\theta :
\mathbbm{R}_+ \rightarrow \mathbbm{R}_+$ to be a smooth function such $\theta
(x) = 1$ if $x \leqslant 1 / 4$ and $\theta (x) = 0$ if $x \geqslant 1 / 2$.
Set $\theta_t = \theta (\langle \mathD \rangle / t)$. Then \
\begin{eqnarray*}
  \theta_t \sigma_s (\mathD) & = & \text{$0$ for $s \geqslant t$,}\\
  \theta_t (\xi) & = & \text{$1$ for $\xi \leqslant c t$ \text{for some} $c >
  0$.}
\end{eqnarray*}
By the Mihlin-H{\"o}rmander theorem we deduce that the operator $\theta_t$ is
bounded on $L^p$ for any $1 < p < \infty$. Let $Z_t^{\flat} \assign \theta_t
Z_t$ and $K^{\flat}_t \assign \theta_t K_t$, then \
\[ Z_t^{\flat} = \theta_t Z_t = \int_0^t \theta_t \langle \mathD \rangle^{-
   1} \sigma_s (\mathD) u_s \mathd s = \int_0^T \theta_t \langle \mathD
   \rangle^{- 1} \sigma_s (\mathD) u_s \mathd s = \theta_t Z_T, \]
so in particular we have $\| Z_t^{\flat} \|_{L^p} \lesssim \| Z_T \|_{L^p}$
for all $t \leqslant T$. On the other hand the functions $Z_T - Z_T^{\flat}$
and $K_T - K_T^{\flat}$ are spectrally supported outside of a ball or radius
$c T$ and
\[ \lambda \bint (\mathbbm{W}^2_T \succ Z_T) K_T = \lambda \bint
   (\mathbbm{W}^2_T \succ Z^{\flat}_T) K_T + \lambda \bint (\mathbbm{W}^2_T
   \succ (Z_T - Z^{\flat}_T)) K_T . \]
We estimate the second part with the following:

\begin{lemma}
  For any small $\varepsilon > 0$ there exists $\delta > 0$ such that
  \begin{eqnarray*}
    \left| \lambda \bint (\mathbbm{W}^2_T \succ (Z_T - Z^{\flat}_T)) K_T
    \right| & \leqslant & C (\varepsilon, \delta) E (\lambda) Q_T +
    \varepsilon \| K \|_{H^{1 - \delta}} + \varepsilon \lambda \| Z_T \|_{L^4}
    .
  \end{eqnarray*}
\end{lemma}

\begin{proof}
  Using the spectral support properties of the various terms we observe that
  \[ \| \mathbbm{W}^2_T \|_{\VV^{- 1 + \delta}} \lesssim \| \mathbbm{W}^2_T
     \|_{\VV^{- 1 - \delta}} T^{2 \delta}, \]
  and
  \[ T^{2 \delta} \| Z_T - Z^{\flat}_T \|_{L^2} \lesssim \| Z_T - Z^{\flat}_T
     \|_{H^{2 \delta}} \lesssim \| Z_T - Z^{\flat}_T \|^{\frac{2 \delta}{1 / 2
     - \delta}}_{H^{1 / 2 - \delta}} \| Z_T - Z^{\flat}_T \|^{\frac{1 / 2 - 3
     \delta}{1 / 2 - \delta}}_{L^2} \]
  \[ \lesssim \| Z_T \|^{\frac{2 \delta}{1 / 2 - \delta}}_{H^{1 / 2 - \delta}}
     \| Z_T \|^{\frac{1 / 2 - 3 \delta}{1 / 2 - \delta}}_{L^2}, \]
  where we used also interpolation and the $L^2$ bound $\| Z^{\flat}_T
  \|_{L^2} \lesssim \| Z_T \|_{L^2}$. We observe also that
  \begin{equation}
    \| Z_T \|_{H^{1 / 2 - \delta}} \leqslant \| K_T \|_{H^{1 / 2 - \delta}}
    + \lambda \| \mathbbm{W}_T^{[3]} \|_{H^{1 / 2 - \delta}}
    \label{eq:bound-Z-K}
  \end{equation}
  Therefore we estimate as follows
  \[ \begin{array}{lll}
       &  & \left| \lambda \bint (\mathbbm{W}^2_T \succ (Z_T - Z^{\flat}_T) )
       K_T \right|\\
       & \lesssim & \lambda \| \mathbbm{W}^2_T \|_{\VV^{- 1 + \delta}} \| Z_T
       - Z^{\flat}_T \|_{L^2} \| K_T \|_{H^{1 - \delta}}\\
       & \lesssim & \lambda \| \mathbbm{W}^2_T \|_{\VV^{- 1 - \delta}} \| K_T
       \|^{1 + \frac{2 \delta}{1 / 2 - \delta}}_{H^{1 - \delta}} \| Z_T
       \|^{\frac{1 / 2 - 3 \delta}{1 / 2 - \delta}}_{L^2}\\
       &  & + \lambda^{1 + \frac{2 \delta}{1 / 2 - \delta}} \|
       \mathbbm{W}^2_T \|_{\VV^{- 1 - \delta}} \| \mathbbm{W}_T^{[3]}
       \|^{\frac{2 \delta}{1 / 2 - \delta}}_{H^{1 / 2 - \delta}} \| K_T
       \|_{H^{1 - \delta}} \| Z_T \|^{\frac{1 / 2 - 3 \delta}{1 / 2 -
       \delta}}_{L^2}\\
       & \lesssim & \lambda^{\frac{3 / 2 - \delta}{4 (1 / 2 - \delta)}} \|
       \mathbbm{W}^2_T \|_{\VV^{- 1 - \delta}} \| K_T \|^{\frac{1 / 2 +
       \delta}{1 / 2 - \delta}}_{H^{1 - \delta}} \lambda^{\frac{1 / 2 - 3
       \delta}{4 (1 / 2 - \delta)}} \| Z_T \|^{\frac{1 / 2 - 3 \delta}{1 / 2 -
       \delta}}_{L^2}\\
       &  & + \lambda^{\frac{3 / 2 + 7 \delta}{4 (1 / 2 - \delta)}} \|
       \mathbbm{W}^2_T \|_{\VV^{- 1 - \delta}} \| \mathbbm{W}_T^{[3]}
       \|^{\frac{2 \delta}{1 / 2 - \delta}}_{H^{1 / 2 - \delta}} \| K_T
       \|_{H^{1 - \delta}} \lambda^{\frac{1 / 2 - 3 \delta}{4 (1 / 2 -
       \delta)}} \| Z_T \|^{\frac{1 / 2 - 3 \delta}{1 / 2 - \delta}}_{L^2}\\
       & \lesssim & C (\varepsilon, \delta) \lambda^{\frac{3 / 2 - \delta}{1
       / 2 - 3 \delta}} \| \mathbbm{W}^2_T \|^{\frac{4 (1 / 2 - \delta)}{1 / 2
       - 3 \delta}}_{\VV^{- 1 - \delta}} + \varepsilon \| K_T \|^2_{H^{1 -
       \delta}} + \varepsilon \lambda \| Z_T \|^4_{L^2}\\
       &  & + C (\varepsilon, \delta) \lambda^{\frac{3 / 2 + 7 \delta}{1 / 2
       + \delta}} \| \mathbbm{W}^2_T \|^{\frac{4 (1 / 2 - \delta)}{1 / 2 +
       \delta}}_{\VV^{- 1 - \delta}} \| \mathbbm{W}_T^{[3]} \|^{\frac{8
       \delta}{1 / 2 + \delta}}_{H^{1 / 2 - \delta}} + \varepsilon \| K_T
       \|^2_{H^{1 - \delta}} + \varepsilon \lambda \| Z_T \|^4_{L^2}
     \end{array} \]
  
\end{proof}

It remains to absorb the remaining term. Applying Ito's formula we obtain
\[ \lambda \bint (\mathbbm{W}^2_T \succ Z^{\flat}_T) K_T = \lambda \int_0^T
   \bint (\mathbbm{W}^2_t \succ Z^{\flat}_t) \dot{K_t} \mathd t + \lambda
   \int_0^T \bint (\mathbbm{W}^2_t \succ \dot{Z}^{\flat}_t) K_t \mathd t +
   \text{martingale} . \]
The second term in the r.h.s. is harmless thanks to the following Lemma

\begin{lemma}
  For any small $\varepsilon > 0$ there exists $\delta > 0$ such that
  \begin{eqnarray*}
    \left| \lambda \int_0^T \bint (\mathbbm{W}^2_t \succ \dot{Z}^{\flat}_t)
    K_t \mathd t \right| & \leqslant & C (\varepsilon, \delta) E (\lambda) Q_T
    + \varepsilon \| K \|^2_{H^{1 - \delta}} + \varepsilon \lambda \| Z_T
    \|^4_{L^4} .
  \end{eqnarray*}
\end{lemma}

\begin{proof}
  First note that $\dot{\theta}_t (n) = (\langle n \rangle / t^2) \dot{\theta}
  (\langle n \rangle / t)$. In particular $Z^{\flat}_t$ is spectrally
  supported in an annulus with inner radius $t / 4$ and outer radius $t / 2$.
  \[ \| \dot{Z}^{\flat}_t \|^2_{L^2} = \sum_{n \in \mathbbm{Z}^3} \left|
     \dot{\theta} \left( \frac{\langle n \rangle}{t} \right) \frac{\langle n
     \rangle}{t^2} \hat{Z}_T (n) \right|^2 \lesssim \sum_{n \in \mathbbm{Z}^3}
     \left| \dot{\theta} \left( \frac{\langle n \rangle}{t} \right)
     \frac{\langle n \rangle}{t} \frac{\langle n \rangle^{3
     \varepsilon}}{t^{1 + 3 \varepsilon}} \hat{Z}_T (n) \right|^2 \lesssim
     \frac{\| Z_T \|^2_{H^{3 \varepsilon}}}{t^{2 + 6 \varepsilon}} . \]
  By Proposition~\ref{paraproductestimate},
  \[ \begin{array}{lll}
       \left| \lambda \int_0^T \bint (\mathbbm{W}^2_t \succ \dot{Z}^{\flat}_t)
       K_t \mathd t \right| & \lesssim & \lambda \int_0^T \| \mathbbm{W}^2_t
       \|_{\VV^{- 1 + \delta}} \| \dot{Z}^{\flat}_t \|_{L^2} \| K_t \|_{H^{1 -
       \delta}} \mathd t\\
       & \lesssim & \lambda \int_0^T \| \mathbbm{W}^2_t \|_{\VV^{- 1 +
       \delta}} \| Z_T \|_{H^{3 \delta}} \| K_t \|_{H^{1 - \delta}}
       \frac{\mathd t}{t^{1 + 3 \delta}}\\
       & \lesssim & \lambda \| Z_T \|_{H^{3 \delta}} \sup_{0 \leqslant t
       \leqslant T} \| K_t \|_{H^{1 - \delta}} \int_0^T \| \mathbbm{W}^2_t
       \|_{\VV^{- 1 - \delta}} \frac{\mathd t}{t^{1 + \delta}}\\
       & \lesssim & \lambda \| Z_T \|^{8 \delta}_{H^{1 / 2 - \delta}} \| Z_T
       \|^{1 - 8 \delta}_{L^2} \sup_{0 \leqslant t \leqslant T} \| K_t
       \|_{H^{1 - \delta}} \int_0^T \| \mathbbm{W}^2_t \|_{\VV^{- 1 - \delta}}
       \frac{\mathd t}{t^{1 + \delta}}\\
       & \lesssim & \lambda \| Z_T \|^{1 - 8 \delta}_{L^2} (\sup_{0 \leqslant
       t \leqslant T} \| K_t \|_{H^{1 - \delta}})^{1 + 8 \delta} \left(
       \int_0^T \| \mathbbm{W}^2_t \|_{\VV^{- 1 - \delta}} \frac{\mathd
       t}{t^{1 + \delta}} \right)\\
       &  & + \lambda^{1 + 8 \delta} \| Z_T \|^{1 - 8 \delta}_{L^2} (\sup_{0
       \leqslant t \leqslant T} \| K_t \|_{H^{1 - \delta}}) \times\\
       &  & \qquad \times \left( \| \mathbbm{W}_T^{[3]} \|^{8 \delta}_{H^{1 /
       2 - \delta}} \int_0^T \| \mathbbm{W}^2_t \|_{\VV^{- 1 - \delta}}
       \frac{\mathd t}{t^{1 + \delta}} \right)\\
       & \lesssim & \varepsilon \lambda \| Z_T \|^4_{L^2} + \varepsilon
       (\sup_{0 \leqslant t \leqslant T} \| K_t \|_{H^{1 - \delta}})^2\\
       &  & + C (\varepsilon, \delta) \lambda^{\frac{3 + 8 \delta}{1 - 8
       \delta}} \left( \int_0^T \| \mathbbm{W}^2_t \|_{\VV^{- 1 - \delta}}
       \frac{\mathd t}{t^{1 + \delta}} \right)^{12}\\
       &  & + \lambda^{\frac{3 + 40 \delta}{1 + 8 \delta}} \left( \|
       \mathbbm{W}_T^{[3]} \|^{8 \delta}_{H^{1 / 2 - \delta}} \int_0^T \|
       \mathbbm{W}^2_t \|_{\VV^{- 1 - \delta}} \frac{\mathd t}{t^{1 + \delta}}
       \right)\\
       &  & + \varepsilon \lambda \| Z_T \|^4_{L^2} + \varepsilon (\sup_{0
       \leqslant t \leqslant T} \| K_t \|_{H^{1 - \delta}})
     \end{array} \]
\end{proof}

So we are left with the singular term $\int_0^T \bint (\mathbbm{W}^2_t \succ
Z^{\flat}_t) \dot{K_t} \mathd t$. In order to handle it we need to introduce a
further decomposition of the drift via a second paracontrolled Ansatz
\begin{equation}
  w_t = - \lambda J_t (\mathbbm{W}^2_t \succ Z^{\flat}_t) + l_t \quad
  \text{and} \quad \dot{K}_t = - \lambda J^2_t (\mathbbm{W}^2_t \succ
  Z^{\flat}_t) + \dot{L}_t, \label{eq:ansatz-2}
\end{equation}
where $L_t \assign I_t (l)$. Now
\begin{equation}
  \begin{array}{lll}
    \| K_t \|^2_{H^{1 - \delta}} & \lesssim & \lambda^2 \int^t_0 \frac{1}{t^{1
    + \delta}} \| \mathbbm{W}^2_s \|^2_{\VV^{- 1 - \delta}} \mathd s \| Z_T
    \|^2_{L^2} + \int^t_0 \| l_s \|_{L^2}^2 \mathd s\\
    & \lesssim & \lambda^3 \left( \int^t_0 \frac{1}{t^{1 + \delta}} \|
    \mathbbm{W}^2_s \|^2_{\VV^{- 1 - \delta}} \mathd s \right)^2 + \lambda \|
    Z_T \|^4_{L^2} + \int^t_0 \| l_s \|_{L^2}^2 \mathd s,
  \end{array} \label{eq:control-K}
\end{equation}
which implies that quadractic functions of the norm $\| K_t \|_{H^{1 -
\delta}}$ with a small coeffient can always be controlled, uniformly in $[0,
T]$, by the positive term
\[ \lambda \bint Z^4_T + \frac{1}{2} \int^T_0 \| l_s \|_{L^2}^2 \mathd s. \]
Using eq.~(\ref{eq:ansatz-2}) we compute


\[ \lambda \int_0^T \bint (\mathbbm{W}^2_t \succ Z^{\flat}_t) \dot{K_t}
   \mathd t + \frac{1}{2} \int_0^T \| w_t \|^2_{L^2} \mathd t = -
   \frac{\lambda^2}{2} \int_0^T \bint (J_t (\mathbbm{W}^2_t \succ
   Z^{\flat}_t))^2 \mathd t + \frac{1}{2} \int_0^T \| l_t \|^2_{L^2} \mathd t.
\]
We now see how we can renormalize the remaining term with the following Lemma.

\begin{lemma}
  For any small $\varepsilon > 0$ there exists a $\delta > 0$ such that
  \begin{eqnarray*}
    \frac{\lambda^2}{2} \left| \int^T_0 \bint ((J_t \mathbbm{W}^2_t \succ
    Z^{\flat}_t))^2 - (J_t \mathbbm{W}^2_t \circ J_t \mathbbm{W}^2_t)
    (Z^{\flat}_t)^2 \right| & \leqslant & C (\varepsilon, \delta) E (\lambda)
    Q_T + \varepsilon \| K_T \|^2_{H^{1 - \delta}} + \varepsilon \lambda \|
    Z_T \|^4_{L^4} .
  \end{eqnarray*}
\end{lemma}

\begin{proof}
  To prove this first note that $t^{1 / 2} J_t$ satisfies the assumptions of
  Proposition~\ref{paraproductleibniz} with $m = - 1$. So
  \[ \| J_t (\mathbbm{W}^2_t \succ Z^{\flat}_t) - (J_t \mathbbm{W}^2_t) \succ
     Z^{\flat}_t \|_{H^{1 / 2 - 2 \delta}} \lesssim \langle t \rangle^{- 1 /
     2} \| \mathbbm{W}^2_t \|_{\VV^{- 1 - \delta}} \| Z^{\flat}_t  \|_{H^{1 /
     2 - \delta}}, \]
  and by Proposition~\ref{multiplierestimate},
  \[ \| J_t (\mathbbm{W}^2_t \succ Z^{\flat}_t) \|_{H^{- 2 \delta}} \lesssim
     \langle t \rangle^{- 1 / 2 - \delta} \| \mathbbm{W}^2_t \|_{\VV^{- 1 -
     \delta}} \| Z^{\flat}_t  \|_{L^2} . \]
  Summarizing these considerations we have
  \[ \left| \frac{\lambda^2}{2} \int_0^T \bint [(J_t (\mathbbm{W}^2_t \succ
     Z^{\flat}_t))^2 - (J_t (\mathbbm{W}^2_t \succ Z^{\flat}_t)) ((J_t
     \mathbbm{W}^2_t) \succ Z^{\flat}_t)] \mathd t \right| \]
  \[ \lesssim \lambda^2 \sup_{t \leqslant T}  [\| Z^{\flat}_t  \|_{L^2} \|
     Z^{\flat}_t  \|_{H^{1 / 2 - \delta}}] \int_0^T \| \mathbbm{W}^2_t
     \|^2_{\VV^{- 1 - \delta}} \frac{\mathd t}{\langle t \rangle^{1 + \delta}}
     . \]
  Again applying Proposition~\ref{multiplierestimate} and
  Proposition~\ref{paraproductleibniz},
  \[ \left| \frac{\lambda^2}{2} \int_0^T \bint (J_t (\mathbbm{W}^2_t \succ
     Z^{\flat}_t))^2 \mathd t - \frac{\lambda^2}{2} \int_0^T \bint ((J_t
     \mathbbm{W}^2_t) \succ Z^{\flat}_t )^2 \mathd t \right| \]
  \[ \lesssim \lambda^2 \sup_{t \leqslant T}  [\| Z^{\flat}_t  \|_{L^2} \|
     Z^{\flat}_t  \|_{H^{1 / 2 - \delta}}] \int_0^T \| \mathbbm{W}^2_t
     \|^2_{\VV^{- 1 - \varepsilon}} \frac{\mathd t}{\langle t \rangle^{1 +
     \delta}} . \]
  From Proposition~\ref{adjointparaproduct} we have
  \[ \left| \frac{\lambda^2}{2} \bint ((J_t \mathbbm{W}^2_t \succ
     Z^{\flat}_t))^2 - \frac{\lambda^2}{2} \bint (J_t \mathbbm{W}^2_t \succ
     Z^{\flat}_t) \circ J_t \mathbbm{W}^2_t Z^{\flat}_t \right| \lesssim
     \lambda^2 \| J_t \mathbbm{W}^2_t \|_{\VV^{- 2 \delta}}^2  \| Z^{\flat}_t
     \|_{H^{4 \delta}}  \| Z^{\flat}_t \|_{L^2} . \]
  The integrability of this term in time follows from the inequality
  \[ \| J_t \mathbbm{W}^2_t \|_{\VV^{- 2 \varepsilon}}^2 \lesssim \langle t
     \rangle^{- 1 - 2 \delta} \| \mathbbm{W}^2_t \|_{\VV^{- 1 - \delta}}^2 .
  \]
  To prove this last bound, recall that $t^{1 / 2} J_t$ is a Fourier
  multiplier with symbol
  \[ \langle k \rangle^{- 1} (- \hat{\rho}' (\langle k \rangle / t) \langle k
     \rangle / t)^{1 / 2} = \langle k \rangle^{- 1} \eta (\langle k \rangle /
     t), \]
  where $\eta$ is a smooth function supported in an annulus of radius $1$.
  Using this observation and applying Proposition~\ref{multiplierestimate}
  gives the estimate. Finally, applying Proposition~\ref{commutatorestimate}
  we get
  \[ \lambda^2 \| (J_t \mathbbm{W}^2_t \succ Z^{\flat}_t) \circ J_t
     \mathbbm{W}^2_t - (J_t \mathbbm{W}^2_t \circ J_t \mathbbm{W}^2_t)
     (Z^{\flat}_t) \|_{L^2} \lesssim \lambda^2 \| J_t \mathbbm{W}^2_t
     \|_{\VV^{- 2 \delta}}^2  \| Z^\flat_t \|_{H^{\delta \varepsilon}} . \]
  Summarizing we have
  \[ \frac{\lambda^2}{2} \left| \int^T_0 \bint ((J_t \mathbbm{W}^2_t \succ
     Z^{\flat}_t))^2 - (J_t \mathbbm{W}^2_t \circ J_t \mathbbm{W}^2_t)
     (Z^{\flat}_t)^2 \right| \]
  \[ \lesssim \lambda^2 \sup_{t \leqslant T}  [\| Z^{\flat}_t  \|_{L^2} \|
     Z^{\flat}_t  \|_{H^{1 / 2 - \delta}}] \int_0^T \| \mathbbm{W}^2_t
     \|^2_{\VV^{- 1 - \varepsilon}} \frac{\mathd t}{\langle t \rangle^{1 +
     \delta}} \]
  \[ \lesssim \varepsilon \left( \frac{1}{2} \sup_{t \leqslant T} \|
     Z^{\flat}_t  \|_{H^{1 / 2 - \delta}}^2 + \lambda \| Z_T \|_{L^2}^4
     \right) + C (\varepsilon, \delta) \lambda^7 \left( \int_0^T \|
     \mathbbm{W}^2_t \|^2_{\VV^{- 1 - \varepsilon}} \frac{\mathd t}{\langle t
     \rangle^{1 + \delta}} \right)^4 \]
  \[ \lesssim \varepsilon \left( \frac{1}{2} \| Z_T \|_{H^{1 / 2 - \delta}}^2
     + \lambda \| Z_T \|_{L^2}^4 \right) + C (\varepsilon, \delta) \lambda^7
     \int_0^T \| \mathbbm{W}^2_t \|^8_{\VV^{- 1 - \varepsilon}} \frac{\mathd
     t}{\langle t \rangle^{1 + \delta}} . \]
  and using eq.~(\ref{eq:bound-Z-K}) to control $\| Z_T \|_{H^{1 / 2 -
  \delta}}$ in terms of $K_T$ we obtain the claim.
\end{proof}

\subsubsection{Bound for the cubic term}

We are now left to estimate the cubic term which we rewrite as
\[ \lambda \bint W_T Z^3_T = \lambda^4 \bint W_T (\mathbbm{W}_T^{[3]})^3 + 3
   \lambda^3 \bint W_T (\mathbbm{W}_T^{[3]})^2 K_T + 3 \lambda^2 \bint W_T
   (\mathbbm{W}_T^{[3]}) K^2_T + \lambda \bint W_T K^3_T . \]
Using Lemma~\ref{lemma:bound-cubic} we establish that
\[ \lambda \left| \bint W_T K^3_T \right| \leqslant E (\lambda) \| W_T
   \|^K_{W^{- 1 / 2 - \varepsilon, p}} + \delta (\| K_T \|^2_{H^{1 -
   \varepsilon}} + \lambda \| K_T \|^4_{L^4}) . \]
\begin{lemma}
  For any small $\varepsilon > 0$ there exists $\delta > 0$ such that
  \[ \left| 3 \lambda^3 \bint W_T (\mathbbm{W}_T^{[3]})^2 K_T \right| + \left|
     3 \lambda^2 \bint W_T \mathbbm{W}_T^{[3]} K^2_T \right| \leqslant C
     (\varepsilon, \delta) E (\lambda) Q_T + \varepsilon \| K_T \|^2_{H^{1 -
     \delta}} + \varepsilon \lambda \| Z_T \|^4_{L^4} . \]
\end{lemma}

\begin{proof}
  We can write, by Proposition~\ref{paralinearization},
  \[ 3 \lambda^3 \left| \bint W_T (\mathbbm{W}_T^{[3]})^2 K_T \right|
     \leqslant 6 \lambda^3 \left| \bint W_T (\mathbbm{W}_T^{[3]} \succ
     \mathbbm{W}_T^{[3]}) K_T \right| + C \lambda^3 \| W_T \|_{\VV^{- 1 / 2
     - \varepsilon}} \| \mathbbm{W}_T^{[3]} \|^2_{\VV^{1 / 2 - \varepsilon}}
     \| K_T \|_{H^{1 - \varepsilon}} . \]
  Decompose
  \[ W_T (\mathbbm{W}_T^{[3]} \succ \mathbbm{W}_T^{[3]}) = W_T \succ
     (\mathbbm{W}_T^{[3]} \succ \mathbbm{W}_T^{[3]}) + W_T \prec
     (\mathbbm{W}_T^{[3]} \succ \mathbbm{W}_T^{[3]}) + W_T \circ
     (\mathbbm{W}_T^{[3]} \succ \mathbbm{W}_T^{[3]}) . \]
  We can estimate the first two terms by
  \[ 6 \lambda^3 \left| \bint W_T \succ (\mathbbm{W}_T^{[3]} \succ
     \mathbbm{W}_T^{[3]}) K_T \right| \lesssim \lambda^3 \| W_T \|_{\VV^{-
     1 / 2 - \varepsilon}} \| \mathbbm{W}_T^{[3]} \|^2_{L^{\infty}} \| K_T
     \|_{H^{1 - \varepsilon}}, \]
  and
  \[ 6 \lambda^3 \left| \bint W_T \prec (\mathbbm{W}_T^{[3]} \succ
     \mathbbm{W}_T^{[3]}) K_T \right| \lesssim \lambda^3 \| W_T \|_{\VV^{-
     1 / 2 - \varepsilon}} \| \mathbbm{W}_T^{[3]} \|_{L^{\infty}} \|
     \mathbbm{W}_T^{[3]} \|_{\VV^{1 / 2 - \varepsilon}} \| K_T \|_{H^{1 -
     \varepsilon}} . \]
  Young's inequality gives then the appropriate result. For the final term we
  use Proposition~\ref{commutatorestimate} to get
  \[ \begin{array}{lll}
       &  & 6 \lambda^3 \left| \bint W_T \circ (\mathbbm{W}_T^{[3]} \succ
       \mathbbm{W}_T^{[3]}) K_T \right|\\
       & \leqslant & 6 \lambda^3 \left| \bint \mathbbm{W}_T^{[3]} (W_T \circ
       \mathbbm{W}_T^{[3]}) K_T \right| + C \lambda^3 \| W_T \|_{\VV^{- 1 /
       2 - \varepsilon}} \| \mathbbm{W}_T^{[3]} \|^2_{\VV^{1 / 2 -
       \varepsilon}} \| K_T \|_{H^{1 - \varepsilon}}\\
       & \lesssim & \lambda^3 \| \mathbbm{W}_T^{[3]} \|_{\VV^{1 / 2 -
       \varepsilon}} \| W_T \circ \mathbbm{W}_T^{[3]} \|_{\VV^{- \varepsilon}}
       \| K_T \|_{H^{1 - \varepsilon}} + \lambda^3 \| W_T \|_{\VV^{- 1 / 2 -
       \varepsilon}} \| \mathbbm{W}_T^{[3]} \|^2_{\VV^{1 / 2 - \varepsilon}}
       \| K_T \|_{H^{1 - \varepsilon}} .\\
       & \lesssim & \lambda^6 C (\varepsilon) \left( \| \mathbbm{W}_T^{[3]}
       \|_{\VV^{1 / 2 - \varepsilon}} \| W_T \circ \mathbbm{W}_T^{[3]}
       \|_{\VV^{- \varepsilon}} + \| W_T \|_{\VV^{- 1 / 2 - \varepsilon}} \|
       \mathbbm{W}_T^{[3]} \|^2_{\VV^{1 / 2 - \varepsilon}} \right)^2 +
       \varepsilon \| K_T \|^2_{H^{1 - \varepsilon}}
     \end{array} \]
  For the last term we estimate
  \begin{eqnarray*}
    \left| \lambda^2 \bint (W_T \mathbbm{W}_T^{[3]}) K^2_T \right| & \lesssim
    & \lambda^2 \| W_T \mathbbm{W}_T^{[3]} \|_{\VV^{- 1 / 2 - \delta}} \| K_T
    \|^2_{H^{1 / 2 + \delta}},\\
    & \lesssim & \lambda^{^{\frac{7 / 2 - 2 \delta}{2 - 2 \delta}}} \| W_T
    \mathbbm{W}_T^{[3]} \|_{\VV^{- 1 / 2 - \delta}} \lambda^{^{\frac{1 / 2 - 2
    \delta}{2 - 2 \delta}}} \| K_T \|^{2 \left( \frac{1 / 2 - 2 \delta}{1 -
    \delta} \right)}_{L^2} \| K_T \|^{2 \left( \frac{1 / 2 + \delta}{1 -
    \delta} \right)}_{H^{1 - \delta}}\\
    & \leqslant & \lambda^{^{\frac{7 / 2 - 2 \delta}{1 / 2 - 6 \delta}}} \|
    W_T \mathbbm{W}_T^{[3]} \|^{\frac{2 - 2 \delta}{1 / 2 - 6 \delta}}_{\VV^{-
    1 / 2 - \delta}} + \lambda \| K_T \|^4_{L^2} + \| K_T \|^2_{H^{1 -
    \delta}}
  \end{eqnarray*}
  and we observe that
  \[ \| W_T \mathbbm{W}_T^{[3]} \|_{\VV^{- 1 / 2 - \delta}} \leqslant \| W_T
     \succ \mathbbm{W}_T^{[3]} \|_{\VV^{- 1 / 2 - \delta}} + \| W_T \circ
     \mathbbm{W}_T^{[3]} \|_{\VV^{- 1 / 2 - \delta}} + \| W_T \prec
     \mathbbm{W}_T^{[3]} \|_{\VV^{- 1 / 2 - \delta}} \]
  \[ \lesssim \| W_T \|_{\VV^{- 1 / 2 - \delta}} \| \mathbbm{W}_T^{[3]}
     \|_{\VV^{1 / 2 - \delta}} + \| W_T \circ \mathbbm{W}_T^{[3]} \|_{\VV^{-
     \varepsilon}} \]
  which allows to conclude using Lemma~\ref{lemma:renormalized-terms} to bound
  $W_T \circ \mathbbm{W}_T^{[3]}$.
\end{proof}

\subsubsection{Renormalization}

Let us summarize the argument so far. Starting from the expression
\[ \Upsilon_T \assign 4 \lambda \bint \llbracket W_T^3 \rrbracket Z_T + 6
   \lambda \bint \llbracket W_T^2 \rrbracket Z_T^2 + 4 \lambda \bint W_T Z_T^3
   - \bint (2 \gamma_T W_T Z_T + \gamma_T Z_T^2) + \lambda \bint  Z_T^4 +
   \int_0^T \| u_s \|^2_{L^2} \mathd s, \]
we performed a first change of variables~(\ref{eq:ansatz-1}) to cancel the
term linear in $Z_T$ and get
\[ \Upsilon_T = \frac{\lambda}{2} \bint \mathbbm{W}^2 Z_T^2 + 4 \lambda \bint
   W_T Z_T^3 - \bint (2 \gamma_T W_T Z_T + \gamma_T Z_T^2) + \lambda \bint 
   Z_T^4 + \frac{1}{2} \int_0^T \| w_s \|^2_{L^2} \mathd s + \cdots \]
where the dots stands for terms which are under control via the positive terms
appearing in the expression. Next, we expanded the term quadratic in $Z_T$ to
get
\[ \Upsilon_T = \lambda^2 \bint \mathbbm{W}^2_T \mathbbm{W}_T^{[3]} K_T +
   \frac{\lambda}{2} \bint \mathbbm{W}^2_T K_T^2 + 4 \lambda \bint W_T Z_T^3 -
   \bint (2 \gamma_T W_T Z_T + \gamma_T Z_T^2) + \lambda \bint  Z_T^4 +
   \frac{1}{2} \int_0^T \| w_s \|^2_{L^2} \mathd s + \cdots \]
Then rewrote it using commutators
\[ \Upsilon_T = \lambda^2 \bint (\mathbbm{W}^2_T \circ \mathbbm{W}_T^{[3]})
   K_T + \lambda \bint (\mathbbm{W}^2_T \succ Z_T) K_T + 4 \lambda \bint W_T
   Z_T^3 \]
\[ - \bint (2 \gamma_T W_T Z_T + \gamma_T Z_T^2) + \lambda \bint  Z_T^4 +
   \frac{1}{2} \int_0^T \| w_s \|^2_{L^2} \mathd s + \cdots \]
Applied an Ito formula to the second term (after a preliminary modification)
to get
\[ \Upsilon_T = \lambda^2 \bint (\mathbbm{W}^2_T \circ \mathbbm{W}_T^{[3]})
   K_T + \lambda \int_0^T \bint (\mathbbm{W}^2_t \succ Z^{\flat}_t) \dot{K_t}
   \mathd t + 4 \lambda \bint W_T Z_T^3 \]
\[ - \bint (2 \gamma_T W_T Z_T + \gamma_T Z_T^2) + \lambda \bint  Z_T^4 +
   \frac{1}{2} \int_0^T \| w_s \|^2_{L^2} \mathd s + \cdots \]
Introduced a second paracontrolled Ansatz~(\ref{eq:ansatz-2}) to cancel the
resulting contributions and obtained
\[ \Upsilon_T = \lambda^2 \bint (\mathbbm{W}^2_T \circ \mathbbm{W}_T^{[3]})
   K_T - \frac{\lambda^2}{2} \int_0^T \bint (J_t (\mathbbm{W}^2_t \succ
   Z^{\flat}_t))^2 \mathd t + 4 \lambda \bint W_T Z_T^3 \]
\[ - \bint (2 \gamma_T W_T Z_T + \gamma_T Z_T^2) + \lambda \bint  Z_T^4 +
   \frac{1}{2} \int_0^T \| l_t \|^2_{L^2} \mathd t + \cdots \]
Next we applied a commutator lemma to the second term to clarify its
renormalizability and expanded the counterterms to get
\[ \Upsilon_T = \lambda^2 \bint (\mathbbm{W}^2_T \circ \mathbbm{W}_T^{[3]})
   K_T - \frac{\lambda^2}{2} \int_0^T \bint (J_t \mathbbm{W}^2_t \circ J_t
   \mathbbm{W}^2_t) (Z^{\flat}_t)^2 \mathd t + 4 \lambda \bint W_T Z_T^3 \]
\[ - \bint (2 \gamma_T W_T K_T + \gamma_T (Z_T^{\flat})^2) + \lambda \bint 
   Z_T^4 + \frac{1}{2} \int_0^T \| l_t \|^2_{L^2} \mathd t + \cdots \]
It now remains to handle the terms
\begin{equation}
  \lambda^2 \bint (\mathbbm{W}^2_T \circ \mathbbm{W}_T^{[3]}) K_T -
  \frac{\lambda^2}{2} \int_0^T \bint (J_t \mathbbm{W}^2_t \circ J_t
  \mathbbm{W}^2_t ) (Z^{\flat}_t)^2 \mathd t, \label{eq:renormalizand-terms}
\end{equation}
which so far are divergent and need to be renormalized. Recall also that we
still have the counterterms we introduced in the beginning, namely
\[ - 2 \gamma_T \bint W_T Z_T - \gamma_T \bint Z^2_T . \]
To handle the first term in eq.~(\ref{eq:renormalizand-terms}) we consider
\[ \lambda^2 \bint (\mathbbm{W}^2_T \circ \mathbbm{W}_T^{[3]}) K_T - 2
   \gamma_T \bint W_T Z_T = \lambda^2 \bint (\mathbbm{W}^2_T \circ
   \mathbbm{W}_T^{[3]} - 2 \gamma_T \lambda^{- 2} W_T) K_T + 2 \lambda
   \gamma_T \bint W_T \mathbbm{W}_T^{[3]} \]
and observe that the average of the second term in the r.h.s. can be absorbed
in $\delta_T$ while for the first we have
\[ \lambda^2 \left| \bint (\mathbbm{W}^2_T \circ \mathbbm{W}_T^{[3]} - 2
   \gamma_T \lambda^{- 2} W_T) K_T \right| \lesssim \lambda^2 \|
   \mathbbm{W}^2_T \circ \mathbbm{W}_T^{[3]} - 2 \lambda^{- 2} \gamma_T W_T
   \|_{\VV^{- 1 / 2 - \varepsilon}} \| K_T \|_{W^{1, 1 / 2 + \varepsilon}} .
\]
Using Lemma~\ref{lemma:renormalized-terms} and eq.~(\ref{eq:control-K}) we
have this term under control. On the other hand, to renormalize the second
term in eq.~(\ref{eq:renormalizand-terms}), we observe that
\begin{equation}
  \gamma_T \bint Z^2_T = \gamma_T \bint (Z^{\flat}_T)^2 + \gamma_T \bint
  Z^{\flat}_T (Z_T - Z^{\flat}_T) + \gamma_T \bint (Z_T - Z^{\flat}_T)^2 .
  \label{eq:decomp-square}
\end{equation}
For the second and third term in eq.~(\ref{eq:decomp-square}) in this last
expression we have the estimates
\[ \left| \gamma_T \bint Z^{\flat}_T (Z_T - Z^{\flat}_T) \right| \leqslant |
   \gamma_T | \| Z^{\flat}_T \|_{L^2} \| Z_T - Z^{\flat}_T \|_{L^2} \lesssim
   \frac{| \gamma_T |}{\langle T \rangle^{1 / 4}} \| Z^{\flat}_T \|_{L^2} \|
   Z_T - Z^{\flat}_T \|_{H^{1 / 4}}, \]
and
\[ \left| \gamma_T \bint (Z_T - Z^{\flat}_T)^2 \right| \leqslant | \gamma_T |
   \| Z_T - Z^{\flat}_T \|^2_{L^2} \lesssim \frac{| \gamma_T |}{\langle T
   \rangle^{1 / 4}} \| Z^{\flat}_T - Z_T \|_{L^2} \| Z_T - Z^{\flat}_T
   \|_{H^{1 / 4}} . \]
These bounds imply that both of them remain bounded provided $\gamma_T$ does
not grow too fast in $T$ which is indeed insured by
Lemma~\ref{lemma:renormalized-terms}. Deriving in $T$ the first term in the
r.h.s. of eq.~(\ref{eq:decomp-square}) we get
\[ \gamma_T \bint (Z^{\flat}_T)^2 = \int_{0}^T \bint \dot{\gamma}_t
   (Z^{\flat}_t)^2 \mathd t + 2 \int^T_{0} \bint \gamma_t Z^{\flat}_t
   \dot{Z}_t^{\flat} \mathd t. \]
For the second term in this expression we can apply the estimate
\[ \left| \int^T_{0} \bint \gamma_t Z^{\flat}_t \dot{Z}_t^{\flat} \mathd t
   \right| \leqslant \int^T_{0} | \gamma_t | \| Z^{\flat}_t \|_{L^2}  \|
   \dot{Z}_t^{\flat} \|_{L^2} \mathd t \lesssim \| Z_T \|_{L^2}  \| Z_T
   \|_{H^{1 / 4}} \int^T_{0} \frac{| \gamma_t | \mathd t}{\langle t
   \rangle^{5 / 4}} . \]
Again, after we have fixed $\gamma_t$ below, we will see this to be bounded.
We conclude that
\begin{eqnarray*}
  \frac{\lambda^2}{2} \int_0^T \bint (J_t \mathbbm{W}^2_t \circ J_t
  \mathbbm{W}^2_t ) (Z^{\flat}_t)^2 \mathd t - \gamma_T \int Z^2_T & = &
  \frac{\lambda^2}{2} \int_0^T \bint (J_t \mathbbm{W}^2_t \circ J_t
  \mathbbm{W}^2_t - 2 \lambda^{- 2} \dot{\gamma}_t) (Z^{\flat}_t)^2 \mathd t +
  R (Z),
\end{eqnarray*}
with
\[ | R (Z) | \lesssim \left[ \frac{| \gamma_T |}{\langle T \rangle^{1 / 4}} +
   \int^T_{0} \frac{| \gamma_t | \mathd t}{\langle t \rangle^{5 / 4}}
   \right] (\| Z_T \|^2_{L^2} + \| Z_T \|^2_{H^{1 / 2 - \delta}}) \lesssim O
   (\lambda^2) (\| Z_T \|^2_{L^2} + \| Z_T \|^2_{H^{1 / 2 - \delta}}) \]
which is therefore under control. Note that we have used that $(\gamma_t)_t$
is of order $\lambda^2$. Finally we have
\[ \left| \lambda^2 \int_0^T \bint (J_t \mathbbm{W}^2_t \circ J_t
   \mathbbm{W}^2_t - 2 \lambda^{- 2} \dot{\gamma}_t) (Z^{\flat}_t)^2 \mathd t
   \right| \]
\[ \lesssim \lambda^2 \left[ \int_0^T \| J_t \mathbbm{W}^2_t \circ J_t
   \mathbbm{W}^2_t - 2 \lambda^{- 2} \dot{\gamma}_t \|_{\VV^{- \varepsilon}}
   \mathd t \right] \sup_{t \in [0, T]} \| (Z^{\flat}_t)^2 \|_{W^{1,
   \varepsilon}} \]
\[ \lesssim \lambda^2 \left[ \int_0^T \| \lambda^2 J_t \mathbbm{W}^2_t \circ
   J_t \mathbbm{W}^2_t - 2 \lambda^{- 2} \dot{\gamma}_t \|_{\VV^{-
   \varepsilon}} \mathd t \right] \| Z_T \|_{H^{\varepsilon}}^2 \]
which can be controlled again using Lemma~\ref{lemma:renormalized-terms}. At
this point we are in the situation where all the terms which are not positive
but depending on the drift can be bounded by
\[ E (\lambda) Q_T + \varepsilon \left( \lambda \bint  Z_T^4 + \frac{1}{2}
   \int_0^T \| l_t \|^2_{L^2} \mathd t \right) \]
with a small constant $\varepsilon$. Therefore
\[ \Upsilon_T \geqslant - E (\lambda) Q_T + (1 - \varepsilon) \left( \lambda
   \bint  Z_T^4 + \frac{1}{2} \int_0^T \| l_t \|^2_{L^2} \mathd t \right)
\]
and taking the expectation we obtain the claimed lower bound in
Theorem~\ref{th:main-3d}.

\subsection{Upper bound}

For the upper bound we note that the bounds we established give
\[ \Upsilon_T \leqslant E (\lambda) Q_T + (1 + \varepsilon) \left( \lambda
   \bint  Z_T^4 + \frac{1}{2} \int_0^T \| l_t \|^2_{L^2} \mathd t \right)
\]
for all admissible drift $u \in \mathbbm{H}_a$. Therefore it is enough to
choose a convenient drift $u$ for which
\begin{equation}
  \sup_T \mathbbm{E} \left( \lambda \bint  Z_T^4 + \frac{1}{2} \int_0^T \|
  l_t \|^2_{L^2} \mathd t \right) < \infty . \label{eq:upper-bound-constraint}
\end{equation}
Choosing $u$ equal to zero, like in the two dimensional case, will not work
since then we would not get a useful lower bound due to the presence of the
diverging constant $\delta_T .$ Instead we need to use a drift corresponding
to the choice $l = 0$ in eq.~(\ref{eq:ansatz-2}), more precisely it would have
to satify the equation
\begin{equation}
  Z_t = - \lambda \mathbbm{W}_t^{[3]} - \int^t_{T_1} J^2_s (\mathbbm{W}^2_s
  \succ Z^{\flat}_s) \mathd s, \label{eq:z-upper-bound}
\end{equation}
for fixed $T_1$ large enough and $t \geqslant T_1$. It will then be sufficient
to find $v$ such that
\begin{equation}
  v_s = - \lambda J_s \mathbbm{W}_s^3 - \lambda \mathbbm{I}_{s \geqslant T_1}
  J_s (\mathbbm{W}^2_s \succ \theta_s I_s (v))  \label{changeofvariables}
\end{equation}
for all $s \geqslant 0$. The reason for introducing $T_1$ is that we can only
solve and control the moments of the solution of eq.~(\ref{changeofvariables})
for large values of $T_1$. We proceed via a fixpoint argument. For any $v \in
\mathbbm{H}_a$ consider the affine map
\[ v \mapsto F (v)_s = - \lambda J_s \mathbbm{W}_s^3 - \lambda \mathbbm{I}_{s
   \geqslant T_1} J_s (\mathbbm{W}^2_s \succ \theta_s I_s (v)) . \]
Now fix some $\alpha < - 1 / 2$, and $\mu > 0$ and introduce the norm
\[ \interleave v \interleave \assign \int_0^{\infty} e^{- \mu s} \| v_s
   \|_{H^{\alpha}}^2 \mathd s. \]
Then observe that
\[ \| F (v)_s \|_{H^{\alpha}} \leqslant \lambda \| J_s \mathbbm{W}_s^3
   \|_{H^{\alpha}} + \lambda \| J_s (\mathbbm{W}^2_s \succ \theta_s I_s (v))
   \|_{H^{\alpha}} \]
\[ \leqslant \lambda \langle s \rangle^{- (1 + \varepsilon) / 2} \|
   \mathbbm{W}_s^3 \|_{H^{\alpha - 1 - \varepsilon}} + \lambda \langle s
   \rangle^{- (1 + \varepsilon) / 2} \| \mathbbm{W}^2_s \|_{\VV^{\alpha - 1 -
   \varepsilon}} \| I_s (v) \|_{L^2} \]
\[ \leqslant C \lambda \langle s \rangle^{- (1 + \varepsilon) / 2} \|
   \mathbbm{W}_s^3 \|_{H^{\alpha - 1 - \varepsilon}} + C \lambda
   \mathbbm{I}_{s \geqslant T_1} \langle s \rangle^{- (1 + \varepsilon) / 2}
   \| \mathbbm{W}^2_s \|_{\VV^{\alpha - 1 - \varepsilon}} \| v \|_{L^2 ([0,
   s], L^2)}, \]
where we used the same computation as in Lemma \ref{lemma:reg-drift}, the fact
that $\theta$ is a bounded operator on $H^{\alpha}$ and the bound on the
operator norm of $\sigma_s (\mathD)$ to get the $\langle s \rangle^{- (1 +
\varepsilon) / 2}$ factor. Therefore we can estimate
\[ \interleave F (v) \interleave = \int_0^{\infty} e^{- \mu s} \| F (v)_s
   \|_{H^{\alpha}}^2 \mathd s \]
\[ \leqslant C \lambda \int_0^{\infty} \| \mathbbm{W}_s^3 \|_{H^{\alpha - 1 -
   \varepsilon}}^2 \frac{e^{- \mu s} \mathd s}{\langle s \rangle^{1 +
   \varepsilon}} + C \lambda \int_0^{\infty} \| v_r \|_{L^2}^2 \mathd r
   \int_r^{\infty} \mathbbm{I}_{s \geqslant T_1} \| \mathbbm{W}^2_s
   \|_{\VV^{\alpha - 1 - \varepsilon}}^2 \frac{e^{- \mu s} \mathd s}{\langle s
   \rangle^{1 + \varepsilon}} \]
\[ \leqslant C \lambda \int_0^{\infty} \| \mathbbm{W}_s^3 \|_{H^{\alpha - 1 -
   \varepsilon}}^2 \frac{e^{- \mu s} \mathd s}{\langle s \rangle^{1 +
   \varepsilon}} \]
\[ \  \]
\[ + C \lambda \int_0^{\infty} \left( \int_r^{\infty} \mathbbm{I}_{s \geqslant
   T_1} \| \mathbbm{W}^2_s \|_{\VV^{\alpha - 1 - \varepsilon}}^4 \frac{\mathd
   s}{\langle s \rangle^{1 + \varepsilon}} \right)^{1 / 2} \left(
   \int_r^{\infty} e^{- 2 \mu s} \mathd s \right)^{1 / 2} \| v_r \|_{L^2}^2
   \mathd r \]
\[ \leqslant C \lambda \int_0^{\infty} \| \mathbbm{W}_s^3 \|_{H^{\alpha - 1 -
   \varepsilon}}^2 \frac{e^{- \mu s} \mathd s}{\langle s \rangle^{1 +
   \varepsilon}} + \frac{C \lambda}{(2 \mu)^{1 / 2}} \left(
   \int_{T_1}^{\infty} \| \mathbbm{W}^2_s \|_{\VV^{\alpha - 1 -
   \varepsilon}}^4 \frac{\mathd s}{\langle s \rangle^{1 + \varepsilon}}
   \right)^{1 / 2} \int_0^{\infty} e^{- \mu r} \| v_r \|_{L^2}^2 \mathd r . \]
\[ \  \]
Choosing $\mu$ large enough so that
\[ \frac{C \lambda}{(2 \mu)^{1 / 2}} \left( \int_{T_1}^{\infty} \|
   \mathbbm{W}^2_s \|_{\VV^{\alpha - 1 - \varepsilon}}^4 \frac{\mathd
   s}{\langle s \rangle^{1 + \varepsilon}} \right)^{1 / 2} \leqslant
   \frac{1}{2}, \]
we have
\[ \interleave F (v) \interleave \leqslant C \lambda \int_0^{\infty} \|
   \mathbbm{W}_s^3 \|_{H^{\alpha - 1 - \varepsilon}}^2 \frac{e^{- \mu s}
   \mathd s}{\langle s \rangle^{1 + \varepsilon}} + \frac{1}{2} \interleave v
   \interleave . \]
Similarly
\[ \interleave F (v) - F (v') \interleave \leqslant \frac{1}{2} \interleave v
   - v' \interleave . \]
These bounds allows to apply a Picard iteration and obtain a drify $v \in
\mathbbm{H}_a$ such that $v = F (v)$ and satisfying the bound
\[ \interleave v \interleave \leqslant 2 C \lambda \int_0^{\infty} \|
   \mathbbm{W}_s^3 \|_{H^{\alpha - 1 - \varepsilon}}^2 \frac{e^{- \mu s}
   \mathd s}{\langle s \rangle^{1 + \varepsilon}} . \]
In order to use this drift in the variational argument we need to control in
expectation the positive term in eq.~(\ref{eq:upper-bound-constraint}). Now
\[ l_s = v_s + \lambda J_s \mathbbm{W}_s^3 + \lambda J_s (\mathbbm{W}^2_s
   \succ \theta_s I_s (v)) = \lambda \mathbbm{I}_{s < T_1} J_s
   (\mathbbm{W}^2_s \succ \theta_s I_s (v)), \qquad s \geqslant 0. \]
Using that $\theta_s I_s (v) = \theta_s I_{T_1} (v)$ for all $0 \leqslant s
\leqslant T_1$ we have
\[ \int_0^T \| l_s \|^2_{L^2} \mathd s \leqslant \lambda \int_0^{T_1} \| J_s
   (\mathbbm{W}^2_s \succ \theta_s I_{T_1} (v)) \|^2_{L^2} \mathd s
   \leqslant \lambda \| I_{T_1} (v) \|_{L^2}^2 \int_0^{T_1} \langle s
   \rangle^{\varepsilon} \| \mathbbm{W}^2_s \|^2_{\VV^{- 1 - \varepsilon}}
   \mathd s \]
\[ \leqslant \lambda \| \lambda \mathbbm{W}_{T_1}^{[3]} \|_{L^2}^4 + \lambda
   \left( \int_0^{T_1} \langle s \rangle^{\varepsilon} \| \mathbbm{W}^2_s
   \|^2_{\VV^{- 1 - \varepsilon}} \mathd s \right)^2, \]
since $I_{T_1} (v) = - \lambda \mathbbm{W}_{T_1}^{[3]}$. The expectation of
this quantity is clearly finite. In order to control $(Z_t)_t$ uniformly in
$t$ we observe that for $t \geqslant T_1$,
\begin{eqnarray*}
  \| Z_t \|_{\VV^{1 / 2 - \varepsilon}} & \leqslant & \lambda \int_0^t \|
  J_s^2 \mathbbm{W}_s^3 \|_{H^{1 / 2 - \varepsilon}} \mathd s + \lambda
  \int^t_{T_1} \| J^2_s (\mathbbm{W}^2_s \succ Z^{\flat}_s ) \|_{H^{1 / 2 -
  \varepsilon}} \mathd s\\
  & \lesssim & \lambda \int_0^t \frac{1}{\langle s \rangle^{1 + \varepsilon}}
  \| \langle \mathD \rangle^{- 2} \mathbbm{W}_s^3 \|_{\VV^{1 / 2 - 2
  \varepsilon}} \mathd s + \lambda \int^t_{T_1} \frac{1}{\langle s \rangle^{1
  + \varepsilon}} \| \langle \mathD \rangle^{- 2} \mathbbm{W}^2_s \|_{\VV^{1 /
  2 - 2 \varepsilon}}  \| Z_s \|_{H^{\alpha}} \mathd s.
\end{eqnarray*}
Using Gronwall's lemma, taking expectation of the $p$-th moment and apply
Cauchy-Schwarz inequality we obtain
\begin{eqnarray*}
  &  & \mathbbm{E} \left[ \| Z_t \|^p_{\VV^{1 / 2 - \varepsilon}} \right]\\
  & \leqslant & \mathbbm{E} \left[ \left( \lambda \int_0^{\infty} \| \langle
  \mathD \rangle^{- 2} \mathbbm{W}_s^3 \|_{\VV^{1 / 2 - 2 \varepsilon}}
  \frac{\mathd s}{\langle s \rangle^{1 + \varepsilon}} \right)^{2 p} \right]
  \mathbbm{E} \left[ \exp \left( 2 p \lambda \int^{\infty}_{T_1} \| \langle
  \mathD \rangle^{- 2} \mathbbm{W}^2_s \|_{\VV^{1 / 2 - 2 \varepsilon}}
  \frac{\mathd s}{\langle s \rangle^{1 + \varepsilon}} \right) \right]\\
  & \lesssim & \lambda^{2 p} \mathbbm{E} \left[ \int_0^{\infty} \| \langle
  \mathD \rangle^{- 2} \mathbbm{W}_s^3 \|^{2 p}_{\VV^{1 / 2 - 2 \varepsilon}}
  \frac{\mathd s}{\langle s \rangle^{1 + \varepsilon}} \right] \mathbbm{E}
  \left[ \int^{\infty}_{T_1} \frac{\mathd s}{s^{1 + \varepsilon / 2}} \exp
  \left( \frac{2 p \lambda}{\langle T_1 \rangle^{\varepsilon / 2}} \| \langle
  \mathD \rangle^{- 2} \mathbbm{W}^2_s \|_{\VV^{1 / 2 - 2 \varepsilon}}
  \right) \right]\\
  & \lesssim & \lambda^{2 p} \mathbbm{E} \left[ \int_0^{\infty} \|
  \mathbbm{W}_s^3 \|^{2 p}_{\VV^{- 3 / 2 - 2 \varepsilon}} \frac{\mathd
  s}{\langle s \rangle^{1 + \varepsilon}} \right] \sup_{s \geqslant T_1}
  \mathbbm{E} \left[ \exp \left( \frac{2 p \lambda}{\langle T_1
  \rangle^{\varepsilon / 2}} \| \langle \mathD \rangle^{- 2} \mathbbm{W}^2_s
  \|_{\VV^{1 / 2 - 2 \varepsilon}} \right) \right],
\end{eqnarray*}
where the second line was obtained by applying Jensen's inequality on the
finite measure space $([T_1, \infty), \langle s \rangle^{- 1 - \varepsilon /
2} \mathd s)$. The first factor is known to be finite and for the second we
have the following lemma which allows to conclude that the average is finite
uniformly in $t$ and prove the lower bound for Theorem~\ref{th:main-3d}.

\begin{lemma}
  If $T_1$ is large enough only depending on $p \lambda$ we have
  \[ \sup_{s \geqslant T_1} \mathbbm{E} \left[ \exp \left( \frac{2 p
     \lambda}{\langle T_1 \rangle^{\varepsilon / 2}} \| \mathbbm{W}^2_s
     \|_{\VV^{- 1 - \varepsilon}} \right) \right] < \infty . \]
\end{lemma}

\begin{proof}
  It is well known that random variables in the second Gaussian chaos are
  exponentially integrable for small coefficients. However in order to give a
  complete proof and show some other uses of the variational formula we give a
  self contained proof in our particular case. Fix $N > 0$. If we apply
  Theorem~\ref{th:variational} with the bounded random variable $V (Y) = 2 p
  \langle T_1 \rangle^{- \varepsilon / 2} \| \mathbbm{Y}^2_s \|_{\VV^{- 1 -
  \varepsilon}} \wedge N$ we deduce that
  \[ \log \mathbbm{E} \left[ \exp \left( \frac{2 p \lambda}{\langle T_1
     \rangle^{\varepsilon / 2}}  \| \mathbbm{Y}^2_s \|_{\VV^{- 1 -
     \varepsilon}} \wedge N \right) \right] \]
  \[ = \sup_{u \in \mathbbm{H}_a} \mathbbm{E} \left[ \frac{2 p
     \lambda}{\langle T_1 \rangle^{\varepsilon / 2}} \| \mathbbm{W}^2_s + 24
     W_s U_s + U^2_s \|_{\VV^{- 1 - \varepsilon}} \wedge N - \frac{| \Lambda
     |}{2} \int^s_0 \| u_t \|^2_{L^2} \mathd t \right] . \]
  \[ \leqslant \sup_{u \in \mathbbm{H}_a} \mathbbm{E} \left[ \frac{2 p
     \lambda}{\langle T_1 \rangle^{\varepsilon / 2}} \| \mathbbm{W}^2_s + 24
     W_s U_s + U^2_s \|_{\VV^{- 1 - \varepsilon}} - \frac{| \Lambda |}{2}
     \int^s_0 \| u_t \|^2_{L^2} \mathd t \right] . \]
  where $U_s = I_s (u)$. By Morrey's inequality and Sobolev embedding, for
  large enough $p$,
  \[ \| U^2_s \|_{\VV^{- 1 - \varepsilon}} \lesssim \| \langle \mathD
     \rangle^{- 1} U^2_s \|_{L^p} \lesssim \| U^2_s \|_{L^3} \lesssim \| U_s
     \|^2_{L^6} \lesssim \| U_s \|^2_{H^1} \lesssim \int^s_0 \| u_t \|^2_{L^2}
     \mathd t. \]
  In order to handle the mixed term we decompose $W_s U_s = W_s \preccurlyeq
  U_s + W_s \succ U_s$ and estimate each term by Besov's embedding and
  paraproduct bounds as
  \[ \| W_s \preccurlyeq U_s \|_{\VV^{- 1 - \varepsilon}} \lesssim \| W_s
     \preccurlyeq U_s \|_{H^{1 / 2 - \varepsilon}} \lesssim \| W_s \|_{\VV^{-
     1 / 2 - \varepsilon}} \| U_s \|_{H^1}, \]
  and
  \[ \| W_s \succ U_s \|_{\VV^{- 1 - \varepsilon}} \lesssim \| W_s \|_{\VV^{-
     1 / 2 - \varepsilon}} \| U_s \|_{\VV^{- 1 / 2}} \lesssim \| W_s
     \|_{\VV^{- 1 / 2 - \varepsilon}} \| U_s \|_{H^1} . \]
  By Young's inequality
  \[ \| W_s U_s \|_{\VV^{- 1 - \varepsilon}} \leqslant \frac{1}{2} \| W_s
     \|_{\VV^{- 1 / 2 - \varepsilon}}^2 + \frac{C}{2} \int^s_0 \| u_t
     \|^2_{L^2} \mathd t. \]
  Therefore, for large enough $T_1 > 0$ (depending only on $p \lambda$)
  \begin{eqnarray*}
    &  & \mathbbm{E} \left[ \frac{2 p \lambda}{\langle T_1
    \rangle^{\varepsilon / 2}} \| (\mathbbm{W}^2_s + 24 W_s U_s + U^2_s)
    \|_{\VV^{- 1 - \varepsilon}} - \frac{| \Lambda |}{2} \int^s_0 \| u_t
    \|^2_{L^2} \mathd t \right]\\
    & \leqslant & \mathbbm{E} \left[ \frac{p \lambda}{\langle T_1
    \rangle^{\varepsilon / 2}} \| \mathbbm{W}^2_s \|_{\VV^{- 1 - \varepsilon}}
    + \frac{(24)^2 p \lambda}{2 \langle T_1 \rangle^{\varepsilon / 2}} \| W_s
    \|_{\VV^{- 1 / 2 - \varepsilon}}^2 + \left[ \frac{Cp \lambda}{\langle T_1
    \rangle^{\varepsilon / 2}} - \frac{| \Lambda |}{2} \right] \int^s_0 \| u_t
    \|^2_{L^2} \mathd t \right]\\
    & \leqslant & \mathbbm{E} \left[ \frac{p \lambda}{\langle T_1
    \rangle^{\varepsilon / 2}} \| \langle \mathD \rangle^{- 2} \mathbbm{W}^2_s
    \|_{\VV^{1 - \varepsilon}} + \frac{(24)^2 p \lambda}{2 \langle T_1
    \rangle^{\varepsilon / 2}} \| W_s \|_{\VV^{- 1 / 2 - \varepsilon}}^2
    \right]
  \end{eqnarray*}
  which is bounded uniformly in $s \geqslant T_1$. This allows us to conclude.
\end{proof}

\section{Estimates of the stochastic terms}

\label{sec:stochastic}In this section we close our argument proving the
following lemmas which give uniform estimates as $T \rightarrow \infty$ of the
stochastic terms appearing in our analytic estimates.

\begin{lemma}
  \label{lemma:renormalized-terms}For any $\varepsilon > 0$ and any $p > 1$
  \begin{equation}
    \sup_T \mathbbm{E} \left[ \| W_T \circ W_T^{[3]} \|_{\VV^{-
    \varepsilon}}^p \right] < \infty . \label{eq:stoch-est-w-w3}
  \end{equation}
  Moreover there exists a function $\gamma_t \in C^1 (\mathbbm{R}_+,
  \mathbbm{R})$ such that for any $\varepsilon > 0$ and any $p > 1$
  \begin{equation}
    \sup_T \mathbbm{E} \left[ \| (\mathbbm{W}^2_T \circ \mathbbm{W}_T^{[3]} -
    2 \gamma_T \lambda^{- 2} W_T) \|_{\VV^{- 1 / 2 - \varepsilon}}^p \right] <
    \infty, \label{eq:renorm-1}
  \end{equation}
  \begin{equation}
    \mathbbm{E} \left[ \left( \int^{\infty}_0 \| J_t \mathbbm{W}^2_t \circ J_t
    \mathbbm{W}^2_t - 2 \lambda^{- 2} \dot{\gamma}_t \|_{\VV^{- \varepsilon}}
    \mathd t \right)^p \right] < \infty . \label{eq:renorm-2}
  \end{equation}
  and
  \begin{equation}
    | \gamma_t | + \langle t \rangle | \dot{\gamma}_t | \lesssim \lambda^2
    \log \langle t \rangle, \qquad t \geqslant 0. \label{eq:bounds-gamma}
  \end{equation}
\end{lemma}

\begin{proof}
  We will concentrate in proving the bounds on the renomalized terms in
  eq.~(\ref{eq:renorm-1}) and (\ref{eq:renorm-2}) and leave to the readed to
  fill the details for the easier term in eq.~(\ref{eq:stoch-est-w-w3}).
  Recall the representation of $W_t = Y_t$ in terms of the family of Brownian
  motions $(B_t^n)_{t, n}$ in eq.~(\ref{eq:rep-Y}). Wick's products of the
  Gaussian field $W_T$ can be represented as iterated stochastic integrals
  wrt. $(B_t^n)_{t, n}$. In particular, if we let $\mathd w_s (k) = \langle k
  \rangle^{- 1} \sigma_s (k) \mathd B_s^k$, we have
  \[ \mathbbm{W}^2_T (x) = 12 \llbracket W_T^2 \rrbracket (x) = 24 \sum_{k_1,
     k_2} e^{i (k_1 + k_2) \cdot x} \int_0^T \int_0^{s_2} \mathd w_{s_1} (k_1)
     \mathd w_{s_2} (k_2), \]
  \[ \mathbbm{W}^{[3]}_T (x) = 24 \sum_{k_1, k_2, k_3} e^{i k_{(123)} \cdot x}
     \int_0^T \int_0^{s_3} \int_0^{s_2} \left( \int_{s_3}^T \frac{\sigma_u^2
     (k_{(123)})}{\langle k_{(123)} \rangle^2} \mathd u \right) \mathd w_{s_1}
     (k_1) \mathd w_{s_2} (k_2) \mathd w_{s_3} (k_3), \]
  where $k_{(123)} \assign k_1 + k_2 + k_3$. Now products of iterated
  integrals can be decomposed in sums of iterated integrals and we get
  \[ \Delta_q (\mathbbm{W}^2_T \circ \mathbbm{W}^{[3]}_T - 2 \gamma_T
     \lambda^{- 2} W_T) (x) = \sum_{k_1, \ldots, k_5} \int_{A_T^5} G^{2 \circ
     [3]}_{0, q} ((s, k)_{1 \cdots 5}) \mathd w_{s_1} (k_1) \cdots \mathd
     w_{s_5} (k_5) \]
  \[ + \sum_{k_1, \ldots, k_3} \int_{A_T^3} G^{2 \circ [3]}_{1, q} ((s, k)_{1
     \cdots 3}) \mathd w_{s_1} (k_1) \cdots \mathd w_{s_3} (k_3) + \sum_{k_1}
     \int_{A_T^1} \tilde{G}^{2 \circ [3]}_{2, q} ((s, k)_1) \mathd w_{s_1}
     (k_1), \]
  where $A^n_T \assign \{ 0 \leqslant s_1 < \cdots < s_n \leqslant T \}
  \subseteq [0, T]^n$ and where the deterministic kernels are given by
  \begin{eqnarray*}
    G^{2 \circ [3]}_{0, q} ((s, k)_{1 \cdots 5}) & = & (24^2) K_q (k_{(1
    \cdots 5)}) e^{i (k_{(1 \cdots 5)}) \cdot x} \sum_{\sigma \in \tmop{Sh}
    (2, 3)} \sum_{i \sim j} \times\\
    &  & \times K_i (k_{(\sigma_1 \sigma_2)}) K_j (k_{(\sigma_3 \sigma_4
    \sigma_5)}) \left( \int_{s_{\sigma_5}}^T \frac{\sigma_u (k_{(\sigma_3
    \sigma_4 \sigma_5)})^2}{\langle k_{(\sigma_3 \sigma_4 \sigma_5)}
    \rangle^2} \mathd u \right),\\
    G^{2 \circ [3]}_{1, q} ((s, k)_{1 \cdots 3}) & = & (24^2) K_q (k_{(1
    \cdots 3)}) e^{i (k_{(1 \cdots 3)}) \cdot x} \sum_{\sigma \in \tmop{Sh}
    (1, 2)} \sum_{i \sim j} \sum_p \int_0^T \mathd r \frac{\sigma_r
    (p)^2}{\langle p \rangle^2} \times\\
    &  & \times K_i (k_{\sigma_1} + q) K_j (k_{(\sigma_2 \sigma_3)} - q)
    \left( \int_{s_{\sigma_3} \vee r}^T \frac{\sigma_u (k_{(\sigma_2
    \sigma_3)} - p)^2}{\langle k_{(\sigma_2 \sigma_3)} - p \rangle^2} \mathd u
    \right),\\
    G^{2 \circ [3]}_{2, q} ((s, k)_1) & = & (24^2) K_q (k_1) e^{i k_1 \cdot x}
    \sum_{i \sim j} \sum_{p_1, p_2} \int_0^T \mathd r_1 \int_0^T \mathd r_2
    \frac{\sigma_{r_1} (p_1)^2}{\langle p_1 \rangle^2} \frac{\sigma_{r_2}
    (p_2)^2}{\langle p_2 \rangle^2} \times\\
    &  & \times K_i (p_1 + p_2) K_j (k_1 - p_1 - p_2) \left( \int_{r_1 \vee
    r_2 \vee s_1}^T \frac{\sigma_u (k_1 - p_1 - p_2)^2}{\langle k_1 - p_1 -
    p_2 \rangle^2} \mathd u \right),\\
    \tilde{G}^{2 \circ [3]}_{2, q} ((s, k)_1) & = & G^{2 \circ [3]}_{2, q}
    ((s, k)_1) - 2 \lambda^{- 2} \gamma_T K_q (k_1) e^{i k_1 \cdot x},
  \end{eqnarray*}
  where $\tmop{Sh} (k, l)$ is the set of permutations of $\{ 1, \ldots, k + l
  \}$ keeping the orders $\sigma (1) < \cdots < \sigma (k)$ and $\sigma (k +
  1) < \cdots < \sigma (k + l)$ and where, for any symbol $z$, we denote with
  expression of the form $z_{1 \cdots n}$ the vector $(z_1, \ldots, z_n)$.
  Estimation of $\Delta_q (\mathbbm{W}^2_T \circ \mathbbm{W}^{[3]}_T) (x)$
  reduces then to estimate each of the three iterated integrals using BDG
  inequalities to get, for any $p \geqslant 2$,
  \[ I_{0, q} = \left\{ \mathbbm{E} \left[ \left| \sum_{k_1, \ldots, k_5}
     \int_{A_T^5} G^{2 \circ [3]}_{0, q} ((s, k)_{1 \cdots 5}) \mathd w_{s_1}
     (k_1) \cdots \mathd w_{s_5} (k_5) \right|^{2 p} \right] \right\}^{1 / p}
  \]
  \[ \lesssim \mathbbm{E} \left[ \left| \sum_{k_1, \ldots, k_5} \int_{A_T^5}
     G^{2 \circ [3]}_{0, q} ((s, k)_{1 \cdots 5}) \mathd w_{s_1} (k_1) \cdots
     \mathd w_{s_5} (k_5) \right|^2 \right] \]
  \[ \lesssim \sum_{k_1, \ldots, k_5} \int_{A_T^5} | G^{2 \circ [3]}_{0, q}
     ((s, k)_{1 \cdots 5}) |^2 \frac{\sigma_{s_1} (k_1)^2}{\langle k_1
     \rangle^2} \cdots \frac{\sigma_{s_5} (k_5)^2}{\langle k_5 \rangle^2}
     \mathd s_1 \cdots \mathd s_5 . \]
  The kernel $G^{2 \circ [3]}_{0, q} ((s, k)_{1 \cdots 5})$ being a symmetric
  function of its argument, we can simplify this expression into an integral
  over $[0, T]^5$:
  \[ \  \]
  \[ I_{0, q} \lesssim \sum_{k_1, \ldots, k_5} \int_{[0, T]^5} | G^{2 \circ
     [3]}_{0, q} ((s, k)_{1 \cdots 5}) |^2 \frac{\sigma_{s_1} (k_1)^2}{\langle
     k_1 \rangle^2} \cdots \frac{\sigma_{s_5} (k_5)^2}{\langle k_5 \rangle^2}
     \mathd s_1 \cdots \mathd s_5 . \]
  Under the measure $\frac{\sigma_{s_5} (k_5)^2}{\langle k_5 \rangle^2} \mathd
  s_5$, we have
  \[ \left| \int_{s_{\sigma_5}}^T \frac{\sigma_u (k_{(\sigma_3 \sigma_4
     \sigma_5)})^2}{\langle k_{(\sigma_3 \sigma_4 \sigma_5)} \rangle^2} \mathd
     u \right| \lesssim \frac{1}{\langle k_{\sigma_5} \rangle^2} . \]
  Therefore with some standard estimates we can reduce us to consider
  \[ I_{0, q} \lesssim \sum_{k_1, \ldots, k_5} \int_{[0, T]^5} \frac{K_q
     (k_{(1 \cdots 5)})^2}{\langle k_5 \rangle^4} \mathbbm{I}_{k_{(12)} \sim
     k_{(345)}} \frac{\sigma_{s_1} (k_1)^2}{\langle k_1 \rangle^2} \cdots
     \frac{\sigma_{s_5} (k_5)^2}{\langle k_5 \rangle^2} \mathd s_1 \cdots
     \mathd s_5 \]
  \[ \lesssim \sum_{k_1, \ldots, k_5} \int_{[0, T]^5} \frac{K_q (k_{(1 \cdots
     5)})^2}{\langle k_5 \rangle^4} \mathbbm{I}_{k_{(12)} \sim k_{(345)}}
     \frac{\sigma_{s_1} (k_1)^2}{\langle k_1 \rangle^2} \cdots
     \frac{\sigma_{s_5} (k_5)^2}{\langle k_5 \rangle^2} \mathd s_1 \cdots
     \mathd s_5 \]
  \[ \lesssim \sum_{k_1, \ldots, k_5} \frac{K_q (k_{(1 \cdots 5)})^2}{\langle
     k_5 \rangle^4} \mathbbm{I}_{k_{(12)} \sim k_{(345)}} \frac{1}{\langle k_1
     \rangle^2} \cdots \frac{1}{\langle k_5 \rangle^2} \]
  \[ \lesssim \sum_{p_1, p_2} \mathbbm{I}_{p_1 \sim p_2} K_q (p_1 + p_2)^2
     \sum_{k_1, \ldots, k_5} \frac{1}{\langle k_5 \rangle^4}
     \mathbbm{I}_{k_{(12)} = p_1, k_{(345)} = p_2} \frac{1}{\langle k_1
     \rangle^2} \cdots \frac{1}{\langle k_5 \rangle^2} \]
  \[ \lesssim \sum_{p_1, p_2} \mathbbm{I}_{p_1 \sim p_2} K_q (p_1 + p_2)^2
     \frac{1}{\langle p_1 \rangle} \frac{1}{\langle p_2 \rangle^4} \lesssim
     \sum_{p_1, r} K_q (r)^2 \frac{1}{\langle p_1 \rangle} \frac{1}{\langle
     p_1 + r \rangle^4} \lesssim \sum_r K_q (r)^2 \frac{1}{\langle r
     \rangle^2} \lesssim 2^q . \]
  Now by similar reasoning we have also
  \[ | G^{2 \circ [3]}_{1, q} ((s, k)_{1 \cdots 3}) | \lesssim \sum_{\sigma
     \in \tmop{Sh} (1, 2)} | K_q (k_{(1 \cdots 3)}) | \sum_{i \sim j} \sum_p
     \int_0^T \mathd r \frac{\sigma_r (p)^2}{\langle p \rangle^2} | K_i
     (k_{\sigma_1} + p) K_j (k_{(\sigma_2 \sigma_3)} - p) | \frac{1}{\langle
     k_{\sigma_1} + p \rangle^2} \]
  \[ \lesssim \sum_{\sigma \in \tmop{Sh} (1, 2)} | K_q (k_{(1 \cdots 3)}) |
     \frac{1}{\langle k_{\sigma_1} \rangle} \]
  so
  \[ I_{1, q} = \left\{ \mathbbm{E} \left[ \left| \sum_{k_1, \ldots, k_3}
     \int_{A_T^3} G^{2 \circ [3]}_{1, q} ((s, k)_{1 \cdots 5}) \mathd y_{s_1}
     (k_1) \cdots \mathd y_{s_3} (k_3) \right|^{2 p} \right] \right\}^{1 / p}
  \]
  \[ \lesssim \sum_{k_1, \ldots, k_3} \int_{[0, T]^5} \left| \sum_{\sigma \in
     \tmop{Sh} (1, 2)} | K_q (k_{(1 \cdots 3)}) | \frac{1}{\langle
     k_{\sigma_1} \rangle} \right|^2 \frac{\sigma_{s_1} (k_1)^2}{\langle k_1
     \rangle^2} \cdots \frac{\sigma_{s_3} (k_3)^2}{\langle k_3 \rangle^2}
     \mathd s_1 \cdots \mathd s_3 \]
  \[ \lesssim \sum_{k_1, \ldots, k_3} | K_q (k_{(1 \cdots 3)}) |^2
     \frac{1}{\langle k_1 \rangle^4} \frac{1}{\langle k_2 \rangle^2}
     \frac{1}{\langle k_3 \rangle^2} \lesssim \sum_r K_q (r)^2
     \frac{1}{\langle r \rangle^2} \lesssim 2^q . \]
  Finally we note that the same strategy cannot be applied to the first chaos,
  since the kernel $G^{2 \circ [3]}_{2, q}$ cannot be uniformly bounded. We
  let
  \[ \begin{array}{lll}
       A_T (s_1, k_1) & = & (24^2) \sum_{i \sim j} \sum_{q_1, q_2} \int_0^T
       \mathd r_1 \int_0^T \mathd r_2 \frac{\sigma_{r_1} (q_1)^2}{\langle q_1
       \rangle^2} \frac{\sigma_{r_2} (q_2)^2}{\langle q_2 \rangle^2} \times\\
       &  & \times K_i (q_1 + q_2) K_j (k_1 - q_1 - q_2) \left( \int_{r_1
       \vee r_2 \vee s_1}^T \frac{\sigma_u^2 (k_1 - q_1 - q_2)}{\langle k_1 -
       q_1 - q_2 \rangle^2} \mathd u \right),
     \end{array} \]
  so
  \[ \tilde{G}^{2 \circ [3]}_{2, q} ((s, k)_1) = K_q (k_1) e^{i k_1 \cdot x}
     [A_T (s_1, k_1) - 2 \lambda^{- 2} \gamma_T] . \]
  Observe that
  \[ \begin{array}{lll}
       A_T (0, 0) & = & (12^2 \cdot 2) \sum_{q_1, q_2} \int_0^T \mathd r_1
       \int_0^T \mathd r_2 \frac{\sigma_{r_1} (q_1)^2}{\langle q_1 \rangle^2}
       \frac{\sigma_{r_2} (q_2)^2}{\langle q_2 \rangle^2} \times\\
       &  & \times \int_{r_1 \vee r_2}^T \frac{\sigma_u^2 (q_1 +
       q_2)}{\langle q_1 + q_2 \rangle^2} \mathd u \sum_{i \sim j} K_i (q_1 +
       q_2) K_j (- q_1 - q_2) .
     \end{array} \]
  We now choose $\gamma_T$ as
  \begin{equation}
    \gamma_T = \lambda^2 A_T (0, 0) = \lambda^2 (12^2 \cdot 2) \sum_{q_1, q_2}
    \int_0^T \mathd u \int_0^u \mathd r_1 \int_0^u \mathd r_2
    \frac{\sigma_{r_1} (q_1)^2}{\langle q_1 \rangle^2} \frac{\sigma_{r_2}
    (q_2)^2}{\langle q_2 \rangle^2} \frac{\sigma_u^2 (q_1 + q_2)}{\langle q_1
    + q_2 \rangle^2} \label{eq:choice-gamma}
  \end{equation}
  where we used the fact that for all $q \in \mathbbm{Z}^d$ we have $\sum_{i
  \sim j} K_i (q) K_j (q) = 1$, since $\int f \circ g = \int f g$. Note that,
  as claimed,
  \[ | \gamma_T | \lesssim \lambda^2 \sum_{q_1, q_2} \frac{\mathbbm{I}_{| q_1
     |, | q_2 |, | q_1 + q_2 | \lesssim T}}{\langle q_1 \rangle^2 \langle q_2
     \rangle^2 \langle q_1 + q_2 \rangle^2} \lesssim \lambda^2 \log \langle T
     \rangle . \]
  Now
  \[ A_T (s_1, k_1) - 2 \lambda^{- 2} \gamma_T = (24^2 \cdot 6) \sum_{q_1,
     q_2} \int_0^T \mathd r_1 \mathd r_2 \frac{\sigma_{r_1} (q_1)^2}{\langle
     q_1 \rangle^2} \frac{\sigma_{r_2} (q_2)^2}{\langle q_2 \rangle^2} \times
  \]
  \[ \times \sum_{i \sim j} K_i (q_1 + q_2) \left( K_j (k_1 - q_1 - q_2)
     \int_{s_1 \vee r_1 \vee r_2}^T \frac{\sigma_u^2 (k_1 - q_1 -
     q_2)}{\langle k_1 - q_1 - q_2 \rangle^2} \mathd u - K_j (q_1 + q_2)
     \int_{r_1 \vee r_2}^T \frac{\sigma_u^2 (q_1 + q_2)}{\langle q_1 + q_2
     \rangle^2} \mathd u \right) \]
  so when $| q_1 + q_2 | \gg | k_1 |$ the quantity in round brackets can be
  estimated by $| k_1 | \langle q_1 + q_2 \rangle^{- 4}$ while when $| q_1 +
  q_2 | \lesssim | k_1 |$ it is estimated by $\langle q_1 + q_2 \rangle^{- 2}$
  so we have
  \[ | A_T (s_1, k_1) - \lambda^{- 2} \gamma_T | \lesssim \sum_{q_1, q_2}
     \frac{1}{\langle q_1 \rangle^2} \frac{1}{\langle q_2 \rangle^2}
     \frac{1}{\langle q_1 + q_2 \rangle^2} \left( \mathbbm{I}_{| q_1 + q_2 |
     \lesssim | k_1 |} +\mathbbm{I}_{| q_1 + q_2 | \gtrsim | k_1 |} \frac{|
     k_1 |}{\langle q_1 + q_2 \rangle^2} \right) \lesssim \log \langle k_1
     \rangle . \]
  And then with this choice of $\gamma_T$ the kernel $\tilde{G}^{2 \circ
  [3]}_{2, q}$ stays uniformly bounded as $T \rightarrow \infty$ and satisfies
  \[ | \tilde{G}^{2 \circ [3]}_{2, q} ((s, k)_1) | \lesssim K_q (k_1) \log
     \langle k_1 \rangle . \]
  From this we easily deduce that
  \[ I_{2, q} = \left\{ \mathbbm{E} \left[ \left| \sum_{k_1} \int_{A_T}
     \tilde{G}^{2 \circ [3]}_{2, q} ((s, k)_{1 \cdots 5}) \mathd y_{s_1} (k_1)
     \right|^{2 p} \right] \right\}^{1 / p} \lesssim q 2^q, \qquad q \geqslant
     - 1. \]
  All together these estimates imply that
  \[ \mathbbm{E} \| \Delta_q (\mathbbm{W}^2_T \circ \mathbbm{W}^{[3]}_T - 2
     \lambda^{- 2} \gamma_T W_T) \|_{L^{2 p}}^{2 p} \lesssim (q 2^{q / 2})^{2
     p}, \qquad q \geqslant - 1. \]
  Standard argument allows to deduce eq.~(\ref{eq:renorm-1}). The analysis of
  the other renormalized product proceeds similarly. Let $V (t) \assign J_t
  \mathbbm{W}^2_t \circ J_t \mathbbm{W}^2_t - 2 \lambda^{- 2} \dot{\gamma}_t$.
  First note that by Besov embedding we have
  \[ \mathbbm{E} \left[ \left( \int^{\infty}_0 \| V (t) \|_{\VV^{-
     \varepsilon}} \mathd t \right)^p \right] \lesssim \mathbbm{E} \left[
     \left( \int^{\infty}_0 \| V (t) \|_{B^{- \varepsilon - d / r}_{r, r}}
     \mathd t \right)^p \right] \]
  \[ \lesssim \mathbbm{E} \left[ \left( \int^{\infty}_0 \left( \sum_q 2^{- q r
     (\varepsilon + d / r)} \| \Delta_q V (t) \|_{L^r} \right)^{1 / r} \mathd
     t \right)^p \right] . \]
  By Minkowski's integral inequality this is bounded by
  \[ \lesssim \left( \int^{\infty}_0 \mathd t \left\{ \mathbbm{E} \left[
     \left( \sum_q 2^{- q r (\varepsilon + d / r)} \| \Delta_q V (t)
     \|_{L^r}^r \right)^{p / r} \right] \right\}^{1 / p} \right)^p . \]
  When $r \geqslant p$ Jensen's inequality and Fubini's theorem give
  \[ \lesssim \left( \int^{\infty}_0 \mathd t \left\{ \sum_q 2^{- q r
     (\varepsilon + d / r)} \int_{\Lambda} \frac{\mathd x}{| \Lambda |}
     \mathbbm{E} [| \Delta_q V (t) (x) |^r] \right\}^{1 / r} \right)^p . \]
  Finally hypercontractivity and stationarity allow to reduce this to bound
  \[ \lesssim \left( \int^{\infty}_0 \mathd t \left\{ \sum_q 2^{- q r
     (\varepsilon + d / r)}  (\mathbbm{E} [| \Delta_q V (t) (0) |^2])^{r /
     2} \right\}^{1 / r} \right)^p . \]
  Letting $I_q (t) =\mathbbm{E} [| \Delta_q V (t) (0) |^2]$ we have
  \[ \mathbbm{E} \left[ \left( \int^{\infty}_0 \| J_t \mathbbm{W}^2_t \circ
     J_t \mathbbm{W}^2_t - 2 \lambda^{- 2} \dot{\gamma}_t \|_{\VV^{-
     \varepsilon}} \mathd t \right)^p \right] \lesssim \left( \int^{\infty}_0
     \mathd t \left\{ \sum_q 2^{- q r (\varepsilon + d / r)}  (I_q (t))^{r /
     2} \right\}^{1 / r} \right)^p . \]
  Now we decompose the random field $\Delta_q (J_t \mathbbm{W}^2_t \circ J_t
  \mathbbm{W}^2_t - 2 \lambda^{- 2} \dot{\gamma}_t) (x)$ into homogeneous
  stochastic integral as above. And obtain
  \[ \Delta_q (\lambda^2 J_t \mathbbm{W}^2_t \circ J_t \mathbbm{W}^2_t - 2
     \lambda^{- 2} \dot{\gamma}_t) (x) = \sum_{k_1, \ldots, k_4} \int_{A_t^4}
     G^{J 2 \circ J 2}_{0, q} ((s, k)_{1 \cdots 4}) \mathd w_{s_1} (k_1)
     \cdots \mathd w_{s_4} (k_4) \]
  \[ + \sum_{k_1, k_2} \int_{A_t^2} G^{J 2 \circ J 2}_{1, q} ((s, k)_{1 2})
     \mathd w_{s_1} (k_1) \mathd w_{s_2} (k_2) + G^{J 2 \circ J 2}_{2, q} \]
  \begin{eqnarray}
    G^{J 2 \circ J 2}_{0, q} ((s, k)_{1 \cdots 4}) & = & (24^2) K_q (k_{(1
    \cdots 4)}) e^{i (k_{(1 \cdots 4)}) \cdot x} \times \nonumber\\
    &  & \times \sum_{\sigma \in \tmop{Sh} (2, 2)} \sum_{i \sim j} K_i
    (k_{(\sigma_1 \sigma_2)}) K_j (k_{(\sigma_3 \sigma_4)}) \frac{\sigma_t
    (k_{(\sigma_1 \sigma_2)})}{\langle k_{(\sigma_1 \sigma_2)} \rangle}
    \frac{\sigma_t (k_{(\sigma_3 \sigma_4)})}{\langle k_{(\sigma_3 \sigma_4)}
    \rangle} \nonumber\\
    G^{J 2 \circ J 2}_{1, q} ((s, k)_{12}) & = & (24^2) K_q (k_{(12)}) e^{i
    (k_{(12)}) \cdot x} \sum_{\sigma \in \tmop{Sh} (1, 1)} \sum_{i \sim j}
    \sum_q \times \nonumber\\
    &  & \times \int_0^t \mathd r \frac{\sigma_r^2 (q)}{\langle q \rangle^2}
    K_i (k_{\sigma_1} + q) K_j (k_{\sigma_2} - q) \left( \frac{\sigma_t
    (k_{\sigma_1} + q)}{\langle k_{\sigma_1} + q \rangle} \frac{\sigma_t
    (k_{\sigma_2} - q)}{\langle k_{\sigma_2} - q \rangle} \right) \nonumber\\
    G^{J 2 \circ J 2}_{2, q} & = & (24^2) \mathbbm{I}_{q = - 1}  \sum_{i \sim
    j} \sum_{q_1, q_2} \int_0^t \mathd r_1 \int_0^t \mathd r_2 \times
    \nonumber\\
    &  & \times \frac{\sigma_{r_1} (q_1)^2}{\langle q_1 \rangle^2}
    \frac{\sigma_{r_2} (q_2)^2}{\langle q_2 \rangle^2} K_i (q_1 + q_2) K_j (-
    q_1 - q_2) \frac{\sigma_t (q_1 + q_2)^2}{\langle q_1 + q_2 \rangle^2}
    \nonumber\\
    &  & - 2 \dot{\gamma}_t \lambda^{- 2} \mathbbm{I}_{q = - 1} . \nonumber
  \end{eqnarray}
  Using our choice of $\gamma_T$ in eq.~(\ref{eq:choice-gamma}) we have that
  \[ \dot{\gamma}_t = \lambda^2 (12^2 \cdot 2) \sum_{q_1, q_2} \int_0^t \mathd
     r_1 \int_0^t \mathd r_2 \frac{\sigma_{r_1} (q_1)^2}{\langle q_1
     \rangle^2} \frac{\sigma_{r_2} (q_2)^2}{\langle q_2 \rangle^2}
     \frac{\sigma_t^2 (q_1 + q_2)}{\langle q_1 + q_2 \rangle^2} \mathd u, \]
  which implies also that
  \[ G^{J 2 \circ J 2}_{2, q} = 0 \]
  and
  \[ | \dot{\gamma}_t | \lesssim \lambda^2 \frac{\log \langle t
     \rangle}{\langle t \rangle}, \]
  as claimed. We pass now to estimate the other two chaoses. The technique is
  the same we used above. Consider first
  \[ I_{0, q} (t) \assign \mathbbm{E} \left[ \left| \sum_{k_1, \ldots, k_4}
     \int_{A_t^4} G^{J 2 \circ J 2}_{0, q} ((s, k)_{1 \cdots 4}) \mathd
     w_{s_1} (k_1) \cdots \mathd w_{s_4} (k_4) \right|^2 \right] \]
  \[ \lesssim \sum_{k_1, \ldots, k_4} \int_{A_t^4} | G^{J 2 \circ J 2}_{0, q}
     ((s, k)_{1 \cdots 4}) |^2 \frac{\sigma_{s_1} (k_1)^2}{\langle k_1
     \rangle^2} \cdots \frac{\sigma_{s_4} (k_4)^2}{\langle k_4 \rangle^2}
     \mathd s_1 \cdots \mathd s_4 \]
  \[ \lesssim \sum_{k_1, \ldots, k_4} \int_{[0, t]^4} | G^{J 2 \circ J 2}_{0,
     q} ((s, k)_{1 \cdots 4}) |^2 \frac{\sigma_{s_1} (k_1)^2}{\langle k_1
     \rangle^2} \cdots \frac{\sigma_{s_4} (k_4)^2}{\langle k_4 \rangle^2}
     \mathd s_1 \cdots \mathd s_4 \]
  \[ \lesssim \sum_{k_1, \ldots, k_4} K_q (k_{(1 \cdots 4)})^2 \int_{[0, t]^4}
     \frac{\sigma_t^2 (k_{(12)})}{\langle k_{(12)} \rangle^2} \frac{\sigma_t^2
     (k_{(34)})}{\langle k_{(34)} \rangle^2} \frac{\sigma_{s_1}
     (k_1)^2}{\langle k_1 \rangle^2} \cdots \frac{\sigma_{s_4}
     (k_4)^2}{\langle k_4 \rangle^2} \mathd s_1 \cdots \mathd s_4 \]
  \[ \lesssim \sum_{k_1, \ldots, k_4} K_q (k_{(1 \cdots 4)})^2
     \frac{\sigma_t^2 (k_{(12)})}{\langle k_{(12)} \rangle^2} \frac{\sigma_t^2
     (k_{(34)})}{\langle k_{(34)} \rangle^2} \frac{1}{\langle k_1 \rangle^2}
     \cdots \frac{1}{\langle k_4 \rangle^2} \]
  \[ \lesssim \frac{\mathbbm{I}_{2^q \lesssim t}}{\langle t \rangle^6}
     \sum_{k_1, \ldots, k_4} K_q (k_{(1 \cdots 4)})^2 \frac{1}{\langle k_1
     \rangle^2} \frac{1}{\langle k_2 \rangle^2} \frac{1}{\langle k_3
     \rangle^2} \frac{1}{\langle k_4 \rangle^2} \lesssim
     \frac{\mathbbm{I}_{2^q \lesssim t}}{\langle t \rangle^6} 2^{4 q} \]
  where we used that $| \sigma_t (x) | \lesssim t^{- 1 / 2} \mathbbm{I}_{x
  \sim t}$. Now taking $\varepsilon + d / r > 0$ we have
  \[ \int^{\infty}_0 \mathd t \left\{ \sum_q 2^{- q r (\varepsilon + d / r)} 
     (I_{0, q} (t))^{r / 2} \right\}^{1 / r} \lesssim \int^{\infty}_0 \mathd t
     \left\{ \sum_{q : 2^q \lesssim t}  \frac{2^{q r (2 - \varepsilon - d /
     r)}}{\langle t \rangle^{3 r}} \right\}^{1 / r} \]
  \[ \lesssim \int^{\infty}_0 \frac{\mathd t}{\langle t \rangle^{1 +
     \varepsilon + d / r}} \lesssim 1. \]
  Taking into account that $| k_1 |, | k_2 | \lesssim t$ we can estimate
  
  
  \[ | G^{J 2 \circ J 2}_{1, q} ((s, k)_{12}) | \lesssim | K_q (k_{(12)}) |
     \sum_p \frac{\mathbbm{I}_{| p | \lesssim t}}{\langle p \rangle^2} \left(
     \frac{\sigma_t (k_1 + p)}{\langle k_1 + p \rangle} \frac{\sigma_t (k_2 -
     p)}{\langle k_2 - p \rangle} \right) \lesssim | K_q (k_{(12)}) | \langle
     t \rangle^{- 2}, \]
  from which we deduce that
  \[ I_{1, q} (t) \assign \mathbbm{E} \left[ \left| \sum_{k_1, k_2}
     \int_{A_t^2} G^{J 2 \circ J 2}_{1, q} ((s, k)_{12}) \mathd w_{s_1} (k_1)
     \mathd w_{s_2} (k_2) \right|^2 \right] \]
  \[ \lesssim \sum_{k_1, k_2} \int_{A_t^2} | G^{J 2 \circ J 2}_{1, q} ((s,
     k)_{12}) |^2 \frac{\sigma_{s_1} (k_1)^2}{\langle k_1 \rangle^2}
     \frac{\sigma_{s_2} (k_2)^2}{\langle k_2 \rangle^2} \mathd s_1 \mathd s_2
  \]
  \[ \lesssim \langle t \rangle^{- 4} \sum_{k_1, k_2} | K_q (k_{(12)}) |^2
     \int_{[0, t]^2} \frac{\sigma_{s_1} (k_1)^2}{\langle k_1 \rangle^2}
     \frac{\sigma_{s_2} (k_2)^2}{\langle k_2 \rangle^2} \mathd s_1 \mathd s_2
  \]
  \[ \lesssim \langle t \rangle^{- 4} \sum_{k_1, k_2} | K_q (k_{(12)}) |^2
     \frac{\mathbbm{I}_{k_1 \lesssim t}}{\langle k_1 \rangle^2}
     \frac{\mathbbm{I}_{k_2 \lesssim t}}{\langle k_2 \rangle^2} \lesssim
     \langle t \rangle^{- 4} 2^{2 q} \mathbbm{I}_{2^q \lesssim t}, \]
  and then, as for $I_{0, q}$, we have
  
  
  \[ \int^{\infty}_0 \mathd t \left\{ \sum_q 2^{- q r (\varepsilon + d / r)} 
     (I_{1, q} (t))^{r / 2} \right\}^{1 / r} \lesssim \int^{\infty}_0
     \frac{\mathd t}{\langle t \rangle^2} \left\{ \sum_q 2^{q r (1 -
     \varepsilon - d / r)} \mathbbm{I}_{2^q \lesssim t} \right\}^{1 / r}
     \lesssim 1, \]
  as claimed. From these estimates standard arguments give
  eq.~(\ref{eq:renorm-2}).
\end{proof}

\appendix\section{Besov spaces and paraproducts}\label{sec:appendix-para}

In this section we will recall some well known results about Besov spaces,
embeddings, Fourier multipliers and paraproducts. The reader can find full
details and proofs
in~{\cite{bahouri_fourier_2011,gubinelli_paracontrolled_2015}}. First recall
the definition of Littlewood--Paley blocks. Let $\chi, \varphi$ be smooth
radial functions $\mathbbm{R}^d \rightarrow \mathbbm{R}$ such that
\begin{itemize}
  \item $\tmop{supp} \chi \subseteq B (0, R)$, $\tmop{supp} \varphi \subseteq
  B (0, 2 R) \setminus B (0, R)$;
  
  \item $0 \leqslant \chi, \varphi \leqslant 1$, $\chi (\xi) + \sum_{j \geq 0}
  \varphi (2^{- j} \xi) = 1$ for any $\xi \in \mathbbm{R}^d$;
  
  \item $\tmop{supp} \varphi (2^{- j} \cdot) \cap \tmop{supp} \varphi (2^{- i}
  \cdot) = \varnothing$ if $| i - j | > 1$.
\end{itemize}
Introduce the notations $\varphi_{- 1} = \chi$, $\varphi_j = \varphi (2^{- j}
\cdot)$ for $j \geqslant 0$. For any $f \in \CS' (\Lambda)$ we define the
operators $\Delta_j f \assign \mathcal{F}^{- 1} \varphi_j (\xi) \hat{f}
(\xi)$, $j \geqslant - 1$.

\begin{definition}
  Let $s \in \mathbbm{R}, p, q \in [1, \infty]$. For a Schwarz distribution $f
  \in \CS' (\Lambda)$ define the norm
  \[ \| f \|_{B_{p, q}^s} = \| (2^{j s} \| \Delta_j f \|_{L^p})_{j \geqslant
     - 1} \|_{\ell^q} \]
  Then the space $B^s_{p, q}$ is the closure of Schwarz distributions under
  this norm. We denote $\VV^{\alpha} = B_{\infty, \infty}^{\alpha}$ the
  Besov--H{\"o}lder space and $H^{\alpha} = B^{\alpha}_{2, 2}$ the Sobolev
  spaces. \ 
\end{definition}

\begin{proposition}
  \label{besovembedding}Let $1 \leqslant p_1 \leqslant p_2 \leqslant \infty, 1
  \leqslant q_1 \leqslant q_2 \leqslant \infty \text{then } B^s_{p_1, q_1}$ is
  continuously included in $B_{p_2, q_2}^{s - d \left( \frac{1}{p_1} -
  \frac{1}{p_2} \right)}$.
\end{proposition}

\begin{proposition}
  \label{compactembedding}For any $s_1, s_2 \in \mathbbm{R}$ such that $s_1 <
  s_2$, any $p, q \in [1, \infty]$ the Besov space $B^{s_1}_{p, q}$ is
  compactly embedded into $B^{s_2}_{p, q}$.
\end{proposition}

\begin{definition}
  Let $ f, g \in \mathcal{\CS} (\Lambda)$. We define the
  paraproducs
  \[ f \succ g = \sum_{j < i - 1} \Delta_i f \Delta_j g, \qquad \text{and}
     \qquad f \prec g = \sum_{j > i + 1} \Delta_i f \Delta_j g = g \succ f. \]
  Moreover we introduce the resonant product
  \[ f \circ g = \sum_{| i - j | \leqslant 1} \Delta_i f \Delta_j g. \]
  Then
  \[ fg = f \prec g + f \circ g + f \succ g. \]
\end{definition}

\begin{proposition}
  \label{paraproductestimate}Let $\alpha < 0, \beta \in \mathbbm{R}$. For
  $ f, g \in \CS (\Lambda)$ we have the estimates
  \begin{eqnarray*}
    \| f \succ g \|_{H^{\beta - \delta}} & \lesssim & \| f \|_{\VV^{\beta}} \|
    g \|_{L^2},\\
    \| f \succ g \|_{\VV^{\beta}} & \lesssim & \| f \|_{\VV^{\beta}} \| g
    \|_{L^{\infty}},\\
    \| f \succ g \|_{H^{\beta - \alpha}} & \lesssim & \| f \|_{\VV^{\beta}} \|
    g \|_{H^{\alpha}},\\
    \| f \succ g \|_{\VV^{\beta}} & \lesssim & \| f \|_{\VV^{\beta}} \| g
    \|_{\VV^{\alpha}} .
  \end{eqnarray*}
  Let $\alpha, \beta \in \mathbbm{R}$ such that $\alpha + \beta > 0$. Then
  \begin{eqnarray*}
    \| f \circ g \|_{H^{\alpha + \beta}} & \lesssim & \| f \|_{\VV^{\beta}} \|
    g \|_{H^{\alpha}},\\
    \| f \circ g \|_{\VV^{\alpha + \beta}} & \lesssim & \| f \|_{\VV^{\beta}}
    \| g \|_{\VV^{\alpha}} .
  \end{eqnarray*}
  By density the paraproduct and resonant product also extend to bilnear
  operators on the respective spaces. 
\end{proposition}

\begin{proposition}
  \label{commutatorestimate}Let $\alpha \in (0, 1)$ $\beta, \gamma \in
  \mathbbm{R}$ such that $\beta + \gamma < 0$, $\alpha + \beta + \gamma > 0$.
  Then for $f, g, h \in \CS$, and for any $\delta > 0$,
  \[ \| (f \succ g) \circ h - g (f \circ h) \|_{H^{\alpha + \beta + \gamma -
     \delta}} \lesssim \| f \|_{\VV^{\gamma}} \| h \|_{\VV^{\beta}} \| g
     \|_{H^{\alpha}}, \]
  \[ \| (f \succ g) \circ h - g (f \circ h) \|_{\VV^{\alpha + \beta + \gamma}}
     \lesssim \| f \|_{\VV^{\gamma}} \| h \|_{\VV^{\beta}} \| g
     \|_{\VV^{\alpha}} . \]
\end{proposition}

\begin{proposition}
  \label{adjointparaproduct}Assume $f \in \VV^{\alpha}, g \in H^{\beta}, h \in
  H^{\gamma}$ and $\alpha + \beta + \gamma = 0$. Then
  \[ \frac{1}{| \mathbbm{T}^d |} \int_{\mathbbm{T}^d} [(f \succ g) h - (f
     \circ h) g] \lesssim \| f \|_{\VV^{\alpha}} \| g \|_{H^{\beta}} \| h
     \|_{H^{\gamma}} . \]
\end{proposition}

\begin{remark}
  Proposition~\ref{adjointparaproduct} is not proven in the above references
  but is quite easy and the reader can fill out a proof.
\end{remark}

\begin{proposition}
  \label{paralinearization}Let $F \in C^2 (\mathbbm{R}, \mathbbm{R})$, and
  assume that $\alpha \in (0, 1)$ then there exists a map $R_F : \VV^{\alpha}
  \rightarrow \VV^{2 \alpha}$ such that for any $f \in \VV^{\alpha}
  (\mathbbm{T}^d)$
  \[ F (f) = f \succ F (f) + R_F (f) . \]
\end{proposition}

\begin{definition}
  A smooth function $\eta$ is said to be an $S^m$ multiplier if for every
  multiindex $\alpha$ there exists a constant $C_{\alpha}$ such that
  \begin{equation}
    \label{symbolinq} \forall \xi \in \mathbbm{R}^d \quad \left|
    \frac{\partial^{\alpha}}{\partial \xi^{\alpha}} f (\xi) \right|
    \lesssim_{\alpha} (1 + | \xi |)^{m - | \alpha |}
  \end{equation}
  We say that a family $\eta_t$ is a uniformly $S^m$ multiplier if
  (\ref{symbolinq}) is satisfied for every $t$ with $C_{\alpha}$ independent
  of $t$. 
\end{definition}

\begin{proposition}
  \label{multiplierestimate}Let $\eta$ be an $S^m$ multiplier, $s \in
  \mathbbm{R}$, $p, q \in [1, \infty]$, and $f \in B_{p, q}^s
  (\mathbbm{T}^d)$, then
  \[ \| \eta (\mathD) f \|_{B_{p, q}^{s - m}} \lesssim \| f \|_{B_{p, q}^s} .
  \]
  Furthermore the constant depends only on $s, p, q, d$ and the constants
  $C_{\alpha}$ in (\ref{symbolinq}).
\end{proposition}

\begin{proposition}
  \label{paraproductleibniz}Assume $m \leqslant 0$, $\alpha \in (0, 1), \beta
  \in \mathbbm{R}$. Let $\eta$ be an $S^m$ multiplier, $f \in \VV^{\beta}$, $g
  \in H^{\alpha}$. Then for any $\delta > 0$.
  \[ \| \eta (\mathD) (f \succ g) - (\eta (\mathD) f \succ g) \|_{H^{\alpha +
     \beta - m - \delta}} \lesssim \| f \|_{\VV^{\beta}} \| g \|_{H^{\alpha}}
  \]
  Again the constant depends only on $\alpha, \beta, \delta$ and the constants
  in (\ref{symbolinq}).
\end{proposition}

\begin{thebibliography}{10}
  \bibitem[1]{albeverio_invariant_2017}Sergio Albeverio  and  Seiichiro
  Kusuoka. {\newblock}The invariant measure and the flow associated to the
  \$\textbackslash Phi {\textasciicircum}4\_3\$-quantum field model.
  {\newblock}Nov 2017.{\newblock}
  
  \bibitem[2]{bahouri_fourier_2011}Hajer Bahouri, Jean-Yves Chemin, and 
  Rapha{\"e}l Danchin. {\newblock}\tmtextit{Fourier Analysis and Nonlinear
  Partial Differential Equations}. {\newblock}Springer, jan 2011.{\newblock}
  
  \bibitem[3]{bauerschmidt_renormalisation_2014}Roland Bauerschmidt,
  David~C.~Brydges, and  Gordon Slade. {\newblock}A renormalisation group
  method. III. Perturbative analysis. {\newblock}\tmtextit{ArXiv:1403.7252
  [math-ph]}, mar 2014. {\newblock}ArXiv: 1403.7252.{\newblock}
  
  \bibitem[4]{benfatto_ultraviolet_1980}G.~Benfatto, M.~Cassandro,
  G.~Gallavotti, F.~Nicol{\'o}, E.~Olivieri, E.~Presutti, and 
  E.~Scacciatelli. {\newblock}Ultraviolet stability in Euclidean scalar field
  theories. {\newblock}\tmtextit{Communications in Mathematical Physics},
  71(2):95--130, jun 1980.{\newblock}
  
  \bibitem[5]{boue_variational_1998}Michelle Bou{\'e}  and  Paul Dupuis.
  {\newblock}A variational representation for certain functionals of Brownian
  motion. {\newblock}\tmtextit{The Annals of Probability}, 26(4):1641--1659,
  oct 1998.{\newblock}
  
  \bibitem[6]{brydges_mayer_1987}D.~C.~Brydges  and  T.~Kennedy.
  {\newblock}Mayer expansions and the Hamilton-Jacobi equation.
  {\newblock}\tmtextit{Journal of Statistical Physics}, 48(1-2):19--49, jul
  1987.{\newblock}
  
  \bibitem[7]{brydges_short_1995}D.~Brydges, J.~Dimock, and  T.~R.~Hurd.
  {\newblock}The short distance behavior of \$(\textbackslash
  phi{\textasciicircum}4)\_3\$. {\newblock}\tmtextit{Communications in
  Mathematical Physics}, 172(1):143--186, 1995.{\newblock}
  
  \bibitem[8]{brydges_lectures_2009}David~C.~Brydges. {\newblock}Lectures on
  the renormalisation group. {\newblock}In \tmtextit{Statistical mechanics}, 
  volume~16  of \tmtextit{IAS/Park City Math. Ser.},  pages  7--93. Amer.
  Math. Soc., Providence, RI, 2009.{\newblock}
  
  \bibitem[9]{brydges_functional_1993}David~C.~Brydges, Roberto Fern{\'a}ndez,
  and  Ch- Ecublens. {\newblock}\tmtextit{Functional Integrals and their
  Applications}. {\newblock}1993.{\newblock}
  
  \bibitem[10]{brydges_new_1983}David~C.~Brydges, J{\"u}rg Fr{\"o}hlich, and 
  Alan~D.~Sokal. {\newblock}A new proof of the existence and nontriviality of
  the continuum \$\textbackslash varphi {\textasciicircum}4\_2\$ and
  \$\textbackslash varphi {\textasciicircum}4\_3\$ quantum field theories.
  {\newblock}\tmtextit{Communications in Mathematical Physics},
  91(2):141--186, 1983.{\newblock}
  
  \bibitem[11]{brydges_renormalisation_2014_3}David~C.~Brydges  and  Gordon
  Slade. {\newblock}A renormalisation group method. I. Gaussian integration
  and normed algebras. {\newblock}\tmtextit{ArXiv:1403.7244 [math-ph]}, mar
  2014. {\newblock}ArXiv: 1403.7244.{\newblock}
  
  \bibitem[12]{brydges_renormalisation_2014_2}David~C.~Brydges  and  Gordon
  Slade. {\newblock}A renormalisation group method. II. Approximation by local
  polynomials. {\newblock}\tmtextit{ArXiv:1403.7253 [math-ph]}, mar 2014.
  {\newblock}ArXiv: 1403.7253.{\newblock}
  
  \bibitem[13]{brydges_renormalisation_2014_1}David~C.~Brydges  and  Gordon
  Slade. {\newblock}A renormalisation group method. IV. Stability analysis.
  {\newblock}\tmtextit{ArXiv:1403.7255 [math-ph]}, mar 2014. {\newblock}ArXiv:
  1403.7255.{\newblock}
  
  \bibitem[14]{brydges_renormalisation_2014}David~C.~Brydges  and  Gordon
  Slade. {\newblock}A renormalisation group method. V. A single
  renormalisation group step. {\newblock}\tmtextit{ArXiv:1403.7256 [math-ph]},
  mar 2014. {\newblock}ArXiv: 1403.7256.{\newblock}
  
  \bibitem[15]{catellier_paracontrolled_2013}R{\'e}mi Catellier  and  Khalil
  Chouk. {\newblock}Paracontrolled Distributions and the 3-dimensional
  Stochastic Quantization Equation. {\newblock}\tmtextit{ArXiv:1310.6869
  [math-ph]}, oct 2013. {\newblock}ArXiv: 1310.6869.{\newblock}
  
  \bibitem[16]{feldman_lambda_1974}Joel Feldman. {\newblock}The
  \$\textbackslash lambda \textbackslash varphi {\textasciicircum}4\_3\$ field
  theory in a finite volume. {\newblock}\tmtextit{Communications in
  Mathematical Physics}, 37:93--120, 1974.{\newblock}
  
  \bibitem[17]{fleming_controlled_2005}Wendell~H.~Fleming  and  Halil~Mete
  Soner. {\newblock}\tmtextit{Controlled Markov Processes and Viscosity
  Solutions}. {\newblock}Springer, 2nd  edition, nov 2005.{\newblock}
  
  \bibitem[18]{follmer_entropy_1985}H.~F{\"o}llmer. {\newblock}An entropy
  approach to the time reversal of diffusion processes. {\newblock}In 
  Prof~M.~Metivier  and  Prof~E.~Pardoux, editors, \tmtextit{Stochastic
  Differential Systems Filtering and Control},  number~69  in  Lecture Notes
  in Control and Information Sciences,  pages  156--163. Springer Berlin
  Heidelberg, 1985.{\newblock}
  
  \bibitem[19]{glimm_boson_1968}James Glimm. {\newblock}Boson fields with the
  \$: \textbackslash Phi{\textasciicircum}\{4\}: \$ interaction in three
  dimensions. {\newblock}\tmtextit{Communications in Mathematical Physics},
  10(1):1--47, 1968.{\newblock}
  
  \bibitem[20]{glimm_positivity_1973}James Glimm  and  Arthur Jaffe.
  {\newblock}Positivity of the ${\varphi}_3^4$ Hamiltonian.
  {\newblock}\tmtextit{Fortschritte der Physik}, 21(7):327--376, jan
  1973.{\newblock}
  
  \bibitem[21]{gubinelli_global_2018}Massimiliano Gubinelli  and  Martina
  Hofmanov{\'a}. {\newblock}Global solutions to elliptic and parabolic
  \$\textbackslash Phi{\textasciicircum}4\$ models in Euclidean space.
  {\newblock}\tmtextit{ArXiv:1804.11253 [math-ph]}, apr 2018.
  {\newblock}ArXiv: 1804.11253.{\newblock}
  
  \bibitem[22]{gubinelli_paracontrolled_2015}Massimiliano Gubinelli, Peter
  Imkeller, and  Nicolas Perkowski. {\newblock}Paracontrolled distributions
  and singular PDEs. {\newblock}\tmtextit{Forum of Mathematics. Pi}, 3:0,
  2015.{\newblock}
  
  \bibitem[23]{gulisashvili_exact_1996}Archil Gulisashvili  and  Mark~A.~Kon.
  {\newblock}Exact Smoothing Properties of Schr{\"o}dinger Semigroups.
  {\newblock}\tmtextit{American Journal of Mathematics}, 118(6):1215--1248,
  1996.{\newblock}
  
  \bibitem[24]{hairer_theory_2014}M.~Hairer. {\newblock}A theory of regularity
  structures. {\newblock}\tmtextit{Inventiones mathematicae}, 198(2):269--504,
  mar 2014.{\newblock}
  
  \bibitem[25]{janson_gaussian_1997}Svante Janson.
  {\newblock}\tmtextit{Gaussian Hilbert Spaces}. {\newblock}Cambridge
  University Press, jun 1997.{\newblock}
  
  \bibitem[26]{kupiainen_renormalization_2016}Antti Kupiainen.
  {\newblock}Renormalization Group and Stochastic PDEs.
  {\newblock}\tmtextit{Annales Henri Poincar{\'e}}, 17(3):497--535, mar
  2016.{\newblock}
  
  \bibitem[27]{lehec_representation_2013}Joseph Lehec.
  {\newblock}Representation formula for the entropy and functional
  inequalities. {\newblock}\tmtextit{Annales de l'Institut Henri Poincar{\'e}
  Probabilit{\'e}s et Statistiques}, 49(3):885--899, 2013.{\newblock}
  
  \bibitem[28]{magnen_infinite_1976}J.~Magnen  and  R.~S{\'e}n{\'e}or.
  {\newblock}The infinite volume limit of the \$\textbackslash phi
  {\textasciicircum}4\_3\$ model. {\newblock}\tmtextit{Ann. Inst. H.
  Poincar{\'e} Sect. A (N.S.)}, 24(2):95--159, 1976.{\newblock}
  
  \bibitem[29]{mourrat_dynamic_2016}Jean-Christophe Mourrat  and  Hendrik
  Weber. {\newblock}The dynamic \$\textbackslash Phi{\textasciicircum}4\_3\$
  model comes down from infinity. {\newblock}\tmtextit{ArXiv:1601.01234
  [math]}, jan 2016. {\newblock}ArXiv: 1601.01234.{\newblock}
  
  \bibitem[30]{otto_quasilinear_2016}Felix Otto  and  Hendrik Weber.
  {\newblock}Quasilinear SPDEs via rough paths.
  {\newblock}\tmtextit{ArXiv:1605.09744 [math]}, may 2016. {\newblock}ArXiv:
  1605.09744.{\newblock}
  
  \bibitem[31]{park_lambda_1975}Yong~Moon Park. {\newblock}The
  \$\textbackslash lambda \textbackslash phi{\textasciicircum}4\_3\$ Euclidean
  quantum field theory in a periodic box. {\newblock}\tmtextit{Journal of
  Mathematical Physics}, 16(11):2183--2188, 1975.{\newblock}
  
  \bibitem[32]{polchinski_renormalization_1984}Joseph Polchinski.
  {\newblock}Renormalization and effective lagrangians.
  {\newblock}\tmtextit{Nuclear Physics B}, 231(2):269--295, jan
  1984.{\newblock}
  
  \bibitem[33]{revuz_continuous_2004}Daniel Revuz  and  Marc Yor.
  {\newblock}\tmtextit{Continuous Martingales and Brownian Motion}.
  {\newblock}Springer, 3rd  edition, dec 2004.{\newblock}
  
  \bibitem[34]{ustunel_variational_2014}Ali~S{\"u}leyman {\"U}st{\"u}nel.
  {\newblock}Variational calculation of Laplace transforms via entropy on
  Wiener space and applications. {\newblock}\tmtextit{Journal of Functional
  Analysis}, 267(8):3058--3083, oct 2014.{\newblock}
  
  \bibitem[35]{wilson_renormalization_1983}Kenneth~G.~Wilson. {\newblock}The
  renormalization group and critical phenomena. {\newblock}\tmtextit{Reviews
  of Modern Physics}, 55(3):583--600, 1983.{\newblock}
  
  \bibitem[36]{zhang_variational_2009}Xicheng Zhang. {\newblock}A variational
  representation for random functionals on abstract Wiener spaces.
  {\newblock}\tmtextit{Journal of Mathematics of Kyoto University},
  49(3):475--490, 2009.{\newblock}
\end{thebibliography}
\end{document}
